\documentclass[8pt]{beamer}
\usetheme{Madrid}
\usecolortheme{default}
\usepackage{amsmath, amssymb, graphicx, array, multirow, booktabs, xcolor, tikz}
\usetikzlibrary{shapes,arrows,positioning,fit,backgrounds}

% Compact font settings
\setbeamerfont{title}{size=\Large}
\setbeamerfont{subtitle}{size=\small}
\setbeamerfont{author}{size=\scriptsize}
\setbeamerfont{date}{size=\scriptsize}
\setbeamerfont{frametitle}{size=\large}
\setbeamerfont{framesubtitle}{size=\normalsize}
\setbeamerfont{enumerate item}{size=\footnotesize}
\setbeamerfont{itemize item}{size=\footnotesize}
\setbeamerfont{block title}{size=\footnotesize}

% Compact spacing
\setlength{\itemsep}{0.08ex}
\setlength{\parskip}{0.08ex}
\setlength{\leftmargini}{0.3cm}
\setbeamersize{text margin left=3mm, text margin right=3mm}
\addtobeamertemplate{frame begin}{\vspace{-0.4cm}}{}

% Colors
\definecolor{myblue}{RGB}{0,0,255}
\definecolor{myred}{RGB}{255,0,0}
\definecolor{mygreen}{RGB}{0,128,0}
\definecolor{myorange}{RGB}{255,165,0}
\definecolor{mypurple}{RGB}{128,0,128}
\definecolor{mygray}{RGB}{128,128,128}
\definecolor{mygold}{RGB}{255,215,0}

\title{Risk and Risk Factors}
\subtitle{100 Questions: Algorithmic Bias, Fairness, Explainability, GenAI Risks}
\author{GARP RAI Exam Preparation}
\date{}

\begin{document}
	
	\begin{frame}
		\titlepage
	\end{frame}
	
	% ============================================================
	% SECTION 1: INTRODUCTION TO AI RISKS - QUESTIONS 1-5
	% ============================================================
	
	\begin{frame}{Q1: AI Risks Overview - Tutorial}
		\fontsize{6.5}{7.5}\selectfont
		\textbf{QUESTION: Why is understanding AI risks particularly challenging?}
		
		\vspace{0.1cm}
		\textbf{OPTIONS:}
		\begin{enumerate}
			\item AI systems are always perfectly reliable
			\item AI is hugely pervasive in range and variety, with rapidly developing capabilities
			\item AI risks are all well-understood and documented
			\item AI systems never fail
		\end{enumerate}
		
		\vspace{0.1cm}
		\textcolor{myblue}{\textbf{ANSWER: B}}
		
		\vspace{0.1cm}
		\textbf{DETAILED EXPLANATION:}
		\begin{itemize}
			\item \textbf{Why B is correct:} AI is used across many domains (finance, healthcare, justice, etc.), and AI capabilities are developing faster than our understanding of risks. This creates challenges for impact prediction, risk measurement, and governance.
			\item \textbf{Why A is wrong:} AI systems are NOT always reliable - they can and do fail
			\item \textbf{Why C is wrong:} Many AI risks are still being discovered and understood
			\item \textbf{Why D is wrong:} AI systems can fail, as demonstrated by numerous real-world examples
		\end{itemize}
		
		\textcolor{myred}{\textbf{EXAM TRAP:}} AI risks are \textcolor{myblue}{evolving, pervasive, and high-stakes} - not static or fully understood!
	\end{frame}
	
	\begin{frame}{Q2: Pervasiveness Challenge - Tutorial}
		\fontsize{6.5}{7.5}\selectfont
		\textbf{QUESTION: Why does the pervasiveness of AI pose a risk management challenge?}
		
		\vspace{0.1cm}
		\textbf{OPTIONS:}
		\begin{enumerate}
			\item AI is only used in one industry
			\item Risks take different forms depending on the use case
			\item All AI risks are identical
			\item Pervasiveness makes risks easier to predict
		\end{enumerate}
		
		\vspace{0.1cm}
		\textcolor{myblue}{\textbf{ANSWER: B}}
		
		\vspace{0.1cm}
		\textbf{DETAILED EXPLANATION:}
		\begin{itemize}
			\item \textbf{Why B is correct:} AI applications span many domains including credit scoring, healthcare, policing, and employment. Each domain has unique risk profiles - a loan denial algorithm has different risks than a medical diagnostic algorithm.
			\item \textbf{Why A is wrong:} AI is used across MANY industries, not just one
			\item \textbf{Why C is wrong:} Risks vary significantly by use case
			\item \textbf{Why D is wrong:} Pervasiveness makes risks HARDER to predict, not easier
		\end{itemize}
		
		\textcolor{myred}{\textbf{EXAM TRAP:}} AI risks are \textcolor{myblue}{context-dependent} across different use cases - no one-size-fits-all approach!
	\end{frame}
	
	\begin{frame}{Q3: High-Stakes AI Use - Tutorial}
		\fontsize{6.5}{7.5}\selectfont
		\textbf{QUESTION: What is a consequence of deploying AI in high-stakes settings?}
		
		\vspace{0.1cm}
		\textbf{OPTIONS:}
		\begin{enumerate}
			\item Technological shortcomings have minimal impact
			\item Technological shortcomings can have significant consequences for people
			\item AI is never used in high-stakes settings
			\item High-stakes use eliminates all risks
		\end{enumerate}
		
		\vspace{0.1cm}
		\textcolor{myblue}{\textbf{ANSWER: B}}
		
		\vspace{0.1cm}
		\textbf{DETAILED EXPLANATION:}
		\begin{itemize}
			\item \textbf{Why B is correct:} Loan applications affect financial well-being, medical diagnoses affect health outcomes, bail decisions affect personal freedom. Algorithmic errors in these domains can have profound impacts on people's lives.
			\item \textbf{Why A is wrong:} High-stakes means impacts are SIGNIFICANT, not minimal
			\item \textbf{Why C is wrong:} AI IS used in high-stakes settings across many domains
			\item \textbf{Why D is wrong:} High-stakes use increases, not eliminates, risk
		\end{itemize}
		
		\textcolor{myred}{\textbf{EXAM TRAP:}} High-stakes AI = \textcolor{myblue}{significant consequences} for individuals - this is why risk management is critical!
	\end{frame}
	
	\begin{frame}{Q4: Rapid AI Development - Tutorial}
		\fontsize{6.5}{7.5}\selectfont
		\textbf{QUESTION: What challenge does rapid AI development create?}
		
		\vspace{0.1cm}
		\textbf{OPTIONS:}
		\begin{enumerate}
			\item It makes risk measurement easier
			\item It creates challenges for impact prediction, risk measurement, and governance
			\item It eliminates all regulatory concerns
			\item It slows down innovation
		\end{enumerate}
		
		\vspace{0.1cm}
		\textcolor{myblue}{\textbf{ANSWER: B}}
		
		\vspace{0.1cm}
		\textbf{DETAILED EXPLANATION:}
		\begin{itemize}
			\item \textbf{Why B is correct:} AI capabilities advance faster than our understanding of risks. Regulatory responses lag behind technology. Governance mechanisms are difficult to establish quickly. This allows potentially harmful practices to persist.
			\item \textbf{Why A is wrong:} Rapid development makes measurement HARDER
			\item \textbf{Why C is wrong:} Regulatory concerns are significant and growing
			\item \textbf{Why D is wrong:} Rapid development ACCELERATES innovation, it doesn't slow it
		\end{itemize}
		
		\textcolor{myred}{\textbf{EXAM TRAP:}} Rapid AI development = \textcolor{myblue}{governance and regulation challenges} - regulation can't keep pace with technology!
	\end{frame}
	
	\begin{frame}{Q5: Interconnected AI Risks - Tutorial}
		\fontsize{6.5}{7.5}\selectfont
		\textbf{QUESTION: How are AI risks typically related?}
		
		\vspace{0.1cm}
		\textbf{OPTIONS:}
		\begin{enumerate}
			\item They are completely independent of each other
			\item They are often interconnected - lack of transparency affects fairness, bias leads to reputational risks
			\item There is only one type of AI risk
			\item Risks are never interconnected
		\end{enumerate}
		
		\vspace{0.1cm}
		\textcolor{myblue}{\textbf{ANSWER: B}}
		
		\vspace{0.1cm}
		\textbf{DETAILED EXPLANATION:}
		\begin{itemize}
			\item \textbf{Why B is correct:} AI risks form a complex network. Lack of transparency makes it impossible to assess fairness. Algorithmic bias leads to reputational damage. Safety failures can reveal fairness issues. Risks reinforce each other.
			\item \textbf{Why A is wrong:} Risks are highly interconnected
			\item \textbf{Why C is wrong:} There are many types of AI risks (functional, data, organizational, strategic, societal)
			\item \textbf{Why D is wrong:} Risks ARE interconnected - they don't exist in isolation
		\end{itemize}
		
		\textcolor{myred}{\textbf{EXAM TRAP:}} AI risks are \textcolor{myblue}{interconnected and mutually reinforcing} - need holistic approach!
	\end{frame}
	
	% ============================================================
	% SECTION 2: ALGORITHMIC BIAS AND FAIRNESS - QUESTIONS 6-30
	% ============================================================
	
	\begin{frame}{Q6: Algorithmic Bias Definition - Tutorial}
		\fontsize{6.5}{7.5}\selectfont
		\textbf{QUESTION: What is algorithmic bias?}
		
		\vspace{0.1cm}
		\textbf{OPTIONS:}
		\begin{enumerate}
			\item Random errors in algorithmic predictions
			\item A systematic deviation of an algorithm's output, performance, or impact relative to some norm or standard
			\item Algorithms that are always perfectly accurate
			\item Errors that only occur in supervised learning
		\end{enumerate}
		
		\vspace{0.1cm}
		\textcolor{myblue}{\textbf{ANSWER: B}}
		
		\vspace{0.1cm}
		\textbf{DETAILED EXPLANATION:}
		\begin{itemize}
			\item \textbf{Why B is correct:} Algorithmic bias is systematic (not random), can affect output, performance, or impact, and is relative to some norm, aim, standard, or baseline. It can be explicit (intentional) or implicit (unintended).
			\item \textbf{Why A is wrong:} Bias is SYSTEMATIC, not random - it follows patterns
			\item \textbf{Why C is wrong:} Algorithms CAN be biased - no algorithm is always accurate
			\item \textbf{Why D is wrong:} Bias can occur in ANY type of learning (supervised, unsupervised, reinforcement)
		\end{itemize}
		
		\textcolor{myred}{\textbf{EXAM TRAP:}} Algorithmic bias = \textcolor{myblue}{systematic deviation from a norm} - not random errors!
	\end{frame}
	
	\begin{frame}{Q7: Explicit vs Implicit Bias - Tutorial}
		\fontsize{6.5}{7.5}\selectfont
		\textbf{QUESTION: What is the difference between explicit and implicit bias?}
		
		\vspace{0.1cm}
		\textbf{OPTIONS:}
		\begin{enumerate}
			\item Explicit bias is always worse than implicit bias
			\item Explicit bias is intentional/projected by developers; implicit bias creeps in undetected
			\item Implicit bias is intentional; explicit bias is accidental
			\item There is no difference
		\end{enumerate}
		
		\vspace{0.1cm}
		\textcolor{myblue}{\textbf{ANSWER: B}}
		
		\vspace{0.1cm}
		\textbf{DETAILED EXPLANATION:}
		\begin{itemize}
			\item \textbf{Why B is correct:} Explicit bias involves the developer intentionally projecting their own biases into the algorithm. Implicit bias creeps in undetected through design choices or training data. Both can lead to unfair outcomes.
			\item \textbf{Why A is wrong:} Both types can be equally harmful to affected individuals
			\item \textbf{Why C is wrong:} This is reversed - explicit is intentional, implicit is unintended
			\item \textbf{Why D is wrong:} They differ in terms of intentionality
		\end{itemize}
		
		\textcolor{myred}{\textbf{EXAM TRAP:}} Explicit = \textcolor{myblue}{intentional}; Implicit = \textcolor{myblue}{unintended} - both can cause harm!
	\end{frame}
	
	\begin{frame}{Q8: Individual Fairness Definition - Tutorial}
		\fontsize{6.5}{7.5}\selectfont
		\textbf{QUESTION: What is individual fairness in algorithmic design?}
		
		\vspace{0.1cm}
		\textbf{OPTIONS:}
		\begin{enumerate}
			\item Treating all groups equally
			\item The doctrine of "treating like cases alike"
			\item Ensuring equal outcomes across demographics
			\item Always using protected characteristics
		\end{enumerate}
		
		\vspace{0.1cm}
		\textcolor{myblue}{\textbf{ANSWER: B}}
		
		\vspace{0.1cm}
		\textbf{DETAILED EXPLANATION:}
		\begin{itemize}
			\item \textbf{Why B is correct:} Individual fairness is based on the Aristotelian doctrine that similarly situated individuals should be treated equally. It focuses on individuals, not groups, and requires agreement on what counts as "alike."
			\item \textbf{Why A is wrong:} This is group fairness, not individual fairness
			\item \textbf{Why C is wrong:} This is demographic parity (a group fairness measure)
			\item \textbf{Why D is wrong:} Individual fairness typically AVOIDS using protected characteristics
		\end{itemize}
		
		\textcolor{myred}{\textbf{EXAM TRAP:}} Individual fairness = \textcolor{myblue}{treating like cases alike} - focus on individuals, not groups!
	\end{frame}
	
	\begin{frame}{Q9: Group Fairness Definition - Tutorial}
		\fontsize{6.5}{7.5}\selectfont
		\textbf{QUESTION: What is group fairness?}
		
		\vspace{0.1cm}
		\textbf{OPTIONS:}
		\begin{enumerate}
			\item Treating each individual uniquely
			\item Examining statistical differences across groups
			\item Only considering individuals with protected characteristics
			\item Ignoring group-level patterns
		\end{enumerate}
		
		\vspace{0.1cm}
		\textcolor{myblue}{\textbf{ANSWER: B}}
		
		\vspace{0.1cm}
		\textbf{DETAILED EXPLANATION:}
		\begin{itemize}
			\item \textbf{Why B is correct:} Group fairness looks at statistical differences across groups sharing protected characteristics. It asks whether certain groups are disadvantaged by the algorithm (e.g., admission rates by gender, loan approvals by race).
			\item \textbf{Why A is wrong:} This is individual fairness
			\item \textbf{Why C is wrong:} Group fairness considers ALL groups, not just those with protected characteristics
			\item \textbf{Why D is wrong:} Group fairness EXAMINES group patterns - that's the whole point
		\end{itemize}
		
		\textcolor{myred}{\textbf{EXAM TRAP:}} Group fairness = \textcolor{myblue}{statistical differences across groups} - looks at group-level patterns!
	\end{frame}
	
	\begin{frame}{Q10: College Admission Example - Tutorial}
		\fontsize{6.5}{7.5}\selectfont
		\textbf{QUESTION: In the college admission example, what fairness issue arises?}
		
		\vspace{0.1cm}
		\textbf{OPTIONS:}
		\begin{enumerate}
			\item All applicants are treated fairly
			\item Individual fairness requires agreeing on what counts as "alike"
			\item Socioeconomic factors never matter
			\item Algorithms always make perfect decisions
		\end{enumerate}
		
		\vspace{0.1cm}
		\textcolor{myblue}{\textbf{ANSWER: B}}
		
		\vspace{0.1cm}
		\textbf{DETAILED EXPLANATION:}
		\begin{itemize}
			\item \textbf{Why B is correct:} What counts as "alike" is often contentious. Should we consider hardships overcome? How to translate resilience into numbers? Different people have different views on what criteria matter.
			\item \textbf{Why A is wrong:} Fairness is not automatically achieved - it requires careful consideration
			\item \textbf{Why C is wrong:} Socioeconomic factors ARE often considered in admission decisions
			\item \textbf{Why D is wrong:} Algorithms can make unfair decisions
		\end{itemize}
		
		\textcolor{myred}{\textbf{EXAM TRAP:}} Individual fairness requires \textcolor{myblue}{agreement on what is "alike"} - which is often contested!
	\end{frame}
	
	\begin{frame}{Q11: Demographic Parity Definition - Tutorial}
		\fontsize{6.5}{7.5}\selectfont
		\textbf{QUESTION: What is demographic parity?}
		
		\vspace{0.1cm}
		\textbf{OPTIONS:}
		\begin{enumerate}
			\item Predictions must be equally accurate across groups
			\item The distribution of predictions must be identical across subpopulations
			\item Positive cases must be correctly identified equally across groups
			\item Precision must be equal across groups
		\end{enumerate}
		
		\vspace{0.1cm}
		\textcolor{myblue}{\textbf{ANSWER: B}}
		
		\vspace{0.1cm}
		\textbf{DETAILED EXPLANATION:}
		\begin{itemize}
			\item \textbf{Why B is correct:} Demographic parity requires $P(\hat{Y}=1|A=0) = P(\hat{Y}=1|A=1)$. The distribution of predictions is identical across groups, ensuring the same outcome rate (e.g., same admission rate across genders).
			\item \textbf{Why A is wrong:} This is predictive rate parity/equal opportunity
			\item \textbf{Why C is wrong:} This is equal opportunity
			\item \textbf{Why D is wrong:} This is predictive rate parity
		\end{itemize}
		
		\textcolor{myred}{\textbf{EXAM TRAP:}} Demographic parity = \textcolor{myblue}{equal distribution of predictions} - independent of correctness!
	\end{frame}
	
	\begin{frame}{Q12: Demographic Parity Formula - Tutorial}
		\fontsize{6.5}{7.5}\selectfont
		\textbf{QUESTION: What is the mathematical definition of demographic parity?}
		
		\vspace{0.1cm}
		\textbf{OPTIONS:}
		\begin{enumerate}
			\item $P(Y=1|\hat{Y}=1,A=0) = P(Y=1|\hat{Y}=1,A=1)$
			\item $P(\hat{Y}=1|A=0) = P(\hat{Y}=1|A=1)$
			\item $P(\hat{Y}=1|Y=1,A=0) = P(\hat{Y}=1|Y=1,A=1)$
			\item $P(Y=1|A=0) = P(Y=1|A=1)$
		\end{enumerate}
		
		\vspace{0.1cm}
		\textcolor{myblue}{\textbf{ANSWER: B}}
		
		\vspace{0.1cm}
		\textbf{DETAILED EXPLANATION:}
		\begin{itemize}
			\item \textbf{Why B is correct:} This is the correct formula where $\hat{Y}$ is the predicted outcome and $A$ is the protected attribute. It states that the probability of a positive prediction is equal regardless of group membership.
			\item \textbf{Why A is wrong:} This is predictive rate parity (uses precision)
			\item \textbf{Why C is wrong:} This is equal opportunity (uses sensitivity)
			\item \textbf{Why D is wrong:} This is base rate equality - a different concept
		\end{itemize}
		
		\textcolor{myred}{\textbf{EXAM TRAP:}} Demographic parity = \textcolor{myblue}{$P(\hat{Y}=1|A=0) = P(\hat{Y}=1|A=1)$} - predictions, not outcomes!
	\end{frame}
	
	\begin{frame}{Q13: Demographic Parity Drawback - Tutorial}
		\fontsize{6.5}{7.5}\selectfont
		\textbf{QUESTION: What is a major drawback of demographic parity?}
		
		\vspace{0.1cm}
		\textbf{OPTIONS:}
		\begin{enumerate}
			\item It is always impossible to achieve
			\item It does not consider the quality/accuracy of predictions
			\item It is too computationally expensive
			\item It requires protected characteristics
		\end{enumerate}
		
		\vspace{0.1cm}
		\textcolor{myblue}{\textbf{ANSWER: B}}
		
		\vspace{0.1cm}
		\textbf{DETAILED EXPLANATION:}
		\begin{itemize}
			\item \textbf{Why B is correct:} Demographic parity is independent of true values. It can be satisfied even if all predictions are wrong (e.g., unqualified candidates labeled as qualified equally across groups). It doesn't ensure predictions are correct.
			\item \textbf{Why A is wrong:} Not always impossible, but can be challenging when base rates differ
			\item \textbf{Why C is wrong:} It's computationally simple - that's part of its appeal
			\item \textbf{Why D is wrong:} It does require knowing group membership, which may not always be available
		\end{itemize}
		
		\textcolor{myred}{\textbf{EXAM TRAP:}} Demographic parity ignores \textcolor{myblue}{prediction accuracy/quality} - a significant limitation!
	\end{frame}
	
	\begin{frame}{Q14: Predictive Rate Parity Definition - Tutorial}
		\fontsize{6.5}{7.5}\selectfont
		\textbf{QUESTION: What is predictive rate parity?}
		
		\vspace{0.1cm}
		\textbf{OPTIONS:}
		\begin{enumerate}
			\item The distribution of predictions is identical across groups
			\item Precision (TP/(TP+FP)) is equal across groups
			\item Sensitivity (TP/(TP+FN)) is equal across groups
			\item Accuracy is equal across groups
		\end{enumerate}
		
		\vspace{0.1cm}
		\textcolor{myblue}{\textbf{ANSWER: B}}
		
		\vspace{0.1cm}
		\textbf{DETAILED EXPLANATION:}
		\begin{itemize}
			\item \textbf{Why B is correct:} Predictive rate parity requires $P(Y=1|\hat{Y}=1,A=0) = P(Y=1|\hat{Y}=1,A=1)$. Among those predicted positive, the proportion correctly identified is the same across groups (equal precision).
			\item \textbf{Why A is wrong:} This is demographic parity
			\item \textbf{Why C is wrong:} This is equal opportunity
			\item \textbf{Why D is wrong:} This is accuracy parity - a different concept
		\end{itemize}
		
		\textcolor{myred}{\textbf{EXAM TRAP:}} Predictive rate parity = \textcolor{myblue}{equal precision across groups} - "of predicted positives, how many were right?"
	\end{frame}
	
	\begin{frame}{Q15: Predictive Rate Parity Drawback - Tutorial}
		\fontsize{6.5}{7.5}\selectfont
		\textbf{QUESTION: What is a drawback of predictive rate parity?}
		
		\vspace{0.1cm}
		\textbf{OPTIONS:}
		\begin{enumerate}
			\item It ignores true positives
			\item It does not account for differences in base rates across groups
			\item It is always impossible to achieve
			\item It requires no data
		\end{enumerate}
		
		\vspace{0.1cm}
		\textcolor{myblue}{\textbf{ANSWER: B}}
		
		\vspace{0.1cm}
		\textbf{DETAILED EXPLANATION:}
		\begin{itemize}
			\item \textbf{Why B is correct:} Base rates may differ significantly across groups. Achieving predictive rate parity becomes extremely challenging when base rates differ, and can have adverse effects on algorithmic performance (may require sacrificing accuracy).
			\item \textbf{Why A is wrong:} It specifically considers true positives through precision
			\item \textbf{Why C is wrong:} Not always impossible, but can be difficult
			\item \textbf{Why D is wrong:} Requires data to calculate precision
		\end{itemize}
		
		\textcolor{myred}{\textbf{EXAM TRAP:}} Predictive rate parity struggles with \textcolor{myblue}{different base rates} - creates trade-offs!
	\end{frame}
	
	\begin{frame}{Q16: Equal Opportunity Definition - Tutorial}
		\fontsize{6.5}{7.5}\selectfont
		\textbf{QUESTION: What is equal opportunity?}
		
		\vspace{0.1cm}
		\textbf{OPTIONS:}
		\begin{enumerate}
			\item The distribution of predictions is identical across groups
			\item Precision is equal across groups
			\item Sensitivity (TP/(TP+FN)) is equal across groups
			\item All predictions are always correct
		\end{enumerate}
		
		\vspace{0.1cm}
		\textcolor{myblue}{\textbf{ANSWER: C}}
		
		\vspace{0.1cm}
		\textbf{DETAILED EXPLANATION:}
		\begin{itemize}
			\item \textbf{Why C is correct:} Equal opportunity requires $P(\hat{Y}=1|Y=1,A=0) = P(\hat{Y}=1|Y=1,A=1)$. The true positive rate (sensitivity) is equal across groups, ensuring positive cases are equally likely to be correctly identified.
			\item \textbf{Why A is wrong:} This is demographic parity
			\item \textbf{Why B is wrong:} This is predictive rate parity
			\item \textbf{Why D is wrong:} No model is always correct - this is impossible
		\end{itemize}
		
		\textcolor{myred}{\textbf{EXAM TRAP:}} Equal opportunity = \textcolor{myblue}{equal true positive rate (sensitivity)} - critical in medical diagnostics!
	\end{frame}
	
	\begin{frame}{Q17: Equal Opportunity Formula - Tutorial}
		\fontsize{6.5}{7.5}\selectfont
		\textbf{QUESTION: What is the mathematical definition of equal opportunity?}
		
		\vspace{0.1cm}
		\textbf{OPTIONS:}
		\begin{enumerate}
			\item $P(\hat{Y}=1|A=0) = P(\hat{Y}=1|A=1)$
			\item $P(Y=1|\hat{Y}=1,A=0) = P(Y=1|\hat{Y}=1,A=1)$
			\item $P(\hat{Y}=1|Y=1,A=0) = P(\hat{Y}=1|Y=1,A=1)$
			\item $P(Y=1|A=0) = P(Y=1|A=1)$
		\end{enumerate}
		
		\vspace{0.1cm}
		\textcolor{myblue}{\textbf{ANSWER: C}}
		
		\vspace{0.1cm}
		\textbf{DETAILED EXPLANATION:}
		\begin{itemize}
			\item \textbf{Why C is correct:} This formula states that the probability of a correct positive prediction ($\hat{Y}=1$) given the actual positive outcome ($Y=1$) is equal across groups (A=0 and A=1).
			\item \textbf{Why A is wrong:} This is demographic parity
			\item \textbf{Why B is wrong:} This is predictive rate parity
			\item \textbf{Why D is wrong:} This is base rate equality
		\end{itemize}
		
		\textcolor{myred}{\textbf{EXAM TRAP:}} Equal opportunity = \textcolor{myblue}{$P(\hat{Y}=1|Y=1,A=0) = P(\hat{Y}=1|Y=1,A=1)$} - equal true positive rate!
	\end{frame}
	
	\begin{frame}{Q18: Equalized Odds Definition - Tutorial}
		\fontsize{6.5}{7.5}\selectfont
		\textbf{QUESTION: What is equalized odds?}
		
		\vspace{0.1cm}
		\textbf{OPTIONS:}
		\begin{enumerate}
			\item Equal distribution of predictions
			\item Equal precision across groups
			\item Equal true positive rate AND equal false positive rate across groups
			\item Equal accuracy across groups
		\end{enumerate}
		
		\vspace{0.1cm}
		\textcolor{myblue}{\textbf{ANSWER: C}}
		
		\vspace{0.1cm}
		\textbf{DETAILED EXPLANATION:}
		\begin{itemize}
			\item \textbf{Why C is correct:} Equalized odds is more restrictive than equal opportunity. It requires both sensitivity equality ($P(\hat{Y}=1|Y=1,A=0) = P(\hat{Y}=1|Y=1,A=1)$) AND specificity equality ($P(\hat{Y}=1|Y=0,A=0) = P(\hat{Y}=1|Y=0,A=1)$).
			\item \textbf{Why A is wrong:} This is demographic parity
			\item \textbf{Why B is wrong:} This is predictive rate parity
			\item \textbf{Why D is wrong:} Accuracy parity is a different concept
		\end{itemize}
		
		\textcolor{myred}{\textbf{EXAM TRAP:}} Equalized odds = \textcolor{myblue}{equal TPR AND equal FPR} - most stringent fairness measure!
	\end{frame}
	
	\begin{frame}{Q19: Fairness Impossibility Theorem - Tutorial}
		\fontsize{6.5}{7.5}\selectfont
		\textbf{QUESTION: What does the fairness impossibility theorem state?}
		
		\vspace{0.1cm}
		\textbf{OPTIONS:}
		\begin{enumerate}
			\item All fairness measures can be satisfied simultaneously
			\item It is impossible to satisfy demographic parity, predictive rate parity, and equal opportunity simultaneously when base rates differ
			\item Fairness is always achievable
			\item Only one fairness measure exists
		\end{enumerate}
		
		\vspace{0.1cm}
		\textcolor{myblue}{\textbf{ANSWER: B}}
		
		\vspace{0.1cm}
		\textbf{DETAILED EXPLANATION:}
		\begin{itemize}
			\item \textbf{Why B is correct:} When base rates differ between populations, it is mathematically impossible to satisfy all three fairness measures (demographic parity, predictive rate parity, and equal opportunity) simultaneously. Trade-offs are unavoidable.
			\item \textbf{Why A is wrong:} This is impossible - the theorem proves it
			\item \textbf{Why C is wrong:} Fairness has inherent trade-offs
			\item \textbf{Why D is wrong:} There are multiple fairness measures
		\end{itemize}
		
		\textcolor{myred}{\textbf{EXAM TRAP:}} Cannot satisfy \textcolor{myblue}{all fairness measures} simultaneously when base rates differ - must prioritize!
	\end{frame}
	
	\begin{frame}{Q20: Accuracy Limitation - Tutorial}
		\fontsize{6.5}{7.5}\selectfont
		\textbf{QUESTION: Why is accuracy not always a good performance measure?}
		
		\vspace{0.1cm}
		\textbf{OPTIONS:}
		\begin{enumerate}
			\item It is always a perfect measure
			\item It can be misleading with unbalanced data
			\item It is too difficult to calculate
			\item It only works for regression
		\end{enumerate}
		
		\vspace{0.1cm}
		\textcolor{myblue}{\textbf{ANSWER: B}}
		
		\vspace{0.1cm}
		\textbf{DETAILED EXPLANATION:}
		\begin{itemize}
			\item \textbf{Why B is correct:} With unbalanced data, high accuracy can be misleading. Example: 100 applicants, 10 suitable, 90 unsuitable - algorithm classifies all as unsuitable, accuracy = 90%, but algorithm provides no useful information.
			\item \textbf{Why A is wrong:} Accuracy has significant limitations
			\item \textbf{Why C is wrong:} Accuracy is easy to calculate
			\item \textbf{Why D is wrong:} Accuracy is for classification, not regression
		\end{itemize}
		
		\textcolor{myred}{\textbf{EXAM TRAP:}} Accuracy is \textcolor{myblue}{misleading with unbalanced data} - always check class distribution!
	\end{frame}
	
	\begin{frame}{Q21: Precision Definition - Tutorial}
		\fontsize{6.5}{7.5}\selectfont
		\textbf{QUESTION: What is precision in classification?}
		
		\vspace{0.1cm}
		\textbf{OPTIONS:}
		\begin{enumerate}
			\item TP/(TP+FN)
			\item TP/(TP+FP)
			\item (TP+TN)/(TP+TN+FP+FN)
			\item TN/(TN+FP)
		\end{enumerate}
		
		\vspace{0.1cm}
		\textcolor{myblue}{\textbf{ANSWER: B}}
		
		\vspace{0.1cm}
		\textbf{DETAILED EXPLANATION:}
		\begin{itemize}
			\item \textbf{Why B is correct:} Precision = TP/(TP+FP). It measures the proportion of positive predictions that are correct. Also called Positive Predictive Value (PPV).
			\item \textbf{Why A is wrong:} This is recall/sensitivity (TP/(TP+FN))
			\item \textbf{Why C is wrong:} This is accuracy
			\item \textbf{Why D is wrong:} This is specificity
		\end{itemize}
		
		\textcolor{myred}{\textbf{EXAM TRAP:}} Precision = \textcolor{myblue}{TP/(TP+FP)} - "of predicted positives, how many were right?"
	\end{frame}
	
	\begin{frame}{Q22: Sensitivity/Recall Definition - Tutorial}
		\fontsize{6.5}{7.5}\selectfont
		\textbf{QUESTION: What is sensitivity (recall) in classification?}
		
		\vspace{0.1cm}
		\textbf{OPTIONS:}
		\begin{enumerate}
			\item TP/(TP+FN)
			\item TP/(TP+FP)
			\item (TP+TN)/(TP+TN+FP+FN)
			\item TN/(TN+FP)
		\end{enumerate}
		
		\vspace{0.1cm}
		\textcolor{myblue}{\textbf{ANSWER: A}}
		
		\vspace{0.1cm}
		\textbf{DETAILED EXPLANATION:}
		\begin{itemize}
			\item \textbf{Why A is correct:} Sensitivity = TP/(TP+FN). It measures the proportion of actual positives that are correctly identified. Also called True Positive Rate (TPR) or Recall.
			\item \textbf{Why B is wrong:} This is precision
			\item \textbf{Why C is wrong:} This is accuracy
			\item \textbf{Why D is wrong:} This is specificity
		\end{itemize}
		
		\textcolor{myred}{\textbf{EXAM TRAP:}} Sensitivity = \textcolor{myblue}{TP/(TP+FN)} - "of actual positives, how many did we catch?"
	\end{frame}
	
	\begin{frame}{Q23: Precision vs Sensitivity Trade-off - Tutorial}
		\fontsize{6.5}{7.5}\selectfont
		\textbf{QUESTION: What trade-off exists between precision and sensitivity?}
		
		\vspace{0.1cm}
		\textbf{OPTIONS:}
		\begin{enumerate}
			\item They always move in the same direction
			\item There is often a trade-off; maximizing one can reduce the other
			\item They are completely independent
			\item Precision is always more important
		\end{enumerate}
		
		\vspace{0.1cm}
		\textcolor{myblue}{\textbf{ANSWER: B}}
		
		\vspace{0.1cm}
		\textbf{DETAILED EXPLANATION:}
		\begin{itemize}
			\item \textbf{Why B is correct:} To catch all positives (high sensitivity), you may increase false positives (lower precision). To avoid false positives (high precision), you may miss positives (lower sensitivity). This is a fundamental trade-off.
			\item \textbf{Why A is wrong:} They often move in opposite directions
			\item \textbf{Why C is wrong:} They are related through the decision threshold
			\item \textbf{Why D is wrong:} Importance depends on context
		\end{itemize}
		
		\textcolor{myred}{\textbf{EXAM TRAP:}} Precision and sensitivity have a \textcolor{myblue}{trade-off} - choosing the right balance depends on context!
	\end{frame}
	
	\begin{frame}{Q24: False Positive Rate Definition - Tutorial}
		\fontsize{6.5}{7.5}\selectfont
		\textbf{QUESTION: What is the False Positive Rate (FPR)?}
		
		\vspace{0.1cm}
		\textbf{OPTIONS:}
		\begin{enumerate}
			\item FP/(TP+FP)
			\item FP/(FP+TN)
			\item FN/(TP+FN)
			\item TN/(TN+FP)
		\end{enumerate}
		
		\vspace{0.1cm}
		\textcolor{myblue}{\textbf{ANSWER: B}}
		
		\vspace{0.1cm}
		\textbf{DETAILED EXPLANATION:}
		\begin{itemize}
			\item \textbf{Why B is correct:} FPR = FP/(FP+TN). It measures the proportion of actual negatives that are incorrectly classified as positive. Also equals 1 - Specificity.
			\item \textbf{Why A is wrong:} This is 1 - precision, not FPR
			\item \textbf{Why C is wrong:} This is False Negative Rate
			\item \textbf{Why D is wrong:} This is specificity
		\end{itemize}
		
		\textcolor{myred}{\textbf{EXAM TRAP:}} FPR = \textcolor{myblue}{FP/(FP+TN)} - important for ROC curves and fairness analysis!
	\end{frame}
	
	\begin{frame}{Q25: Confusion Matrix Components - Tutorial}
		\fontsize{6.5}{7.5}\selectfont
		\textbf{QUESTION: What are the four components of a confusion matrix?}
		
		\vspace{0.1cm}
		\textbf{OPTIONS:}
		\begin{enumerate}
			\item TP, TN, FP, FN
			\item Mean, Median, Mode, Standard Deviation
			\item Training, Validation, Test, Production
			\item Precision, Recall, Accuracy, F1
		\end{enumerate}
		
		\vspace{0.1cm}
		\textcolor{myblue}{\textbf{ANSWER: A}}
		
		\vspace{0.1cm}
		\textbf{DETAILED EXPLANATION:}
		\begin{itemize}
			\item \textbf{Why A is correct:} TP (True Positive), TN (True Negative), FP (False Positive), FN (False Negative) are the four components that form the confusion matrix, which is the foundation for classification metrics.
			\item \textbf{Why B is wrong:} These are statistical measures, not confusion matrix components
			\item \textbf{Why C is wrong:} These are data splits, not confusion matrix components
			\item \textbf{Why D is wrong:} These are performance metrics derived from the confusion matrix
		\end{itemize}
		
		\textcolor{myred}{\textbf{EXAM TRAP:}} Confusion matrix = \textcolor{myblue}{TP, TN, FP, FN} - foundation for all classification metrics!
	\end{frame}
	
	\begin{frame}{Q26: Fairness Measures Summary Table - Tutorial}
		\fontsize{6.5}{7.5}\selectfont
		\textbf{QUESTION: What question does demographic parity answer?}
		
		\vspace{0.1cm}
		\textbf{OPTIONS:}
		\begin{enumerate}
			\item "Is the likelihood of a true positive prediction the same across groups?"
			\item "Is the likelihood of a positive prediction the same across groups?"
			\item "Is the likelihood of identifying actual positives as positive the same across groups?"
			\item "Are all predictions accurate?"
		\end{enumerate}
		
		\vspace{0.1cm}
		\textcolor{myblue}{\textbf{ANSWER: B}}
		
		\vspace{0.1cm}
		\textbf{DETAILED EXPLANATION:}
		\begin{itemize}
			\item \textbf{Why B is correct:} Demographic parity asks whether the probability of a positive prediction ($\hat{Y}=1$) is the same regardless of group membership. It focuses on equal outcome distribution.
			\item \textbf{Why A is wrong:} This is predictive rate parity
			\item \textbf{Why C is wrong:} This is equal opportunity
			\item \textbf{Why D is wrong:} This is accuracy
		\end{itemize}
		
		\textcolor{myred}{\textbf{EXAM TRAP:}} Demographic parity = \textcolor{myblue}{equal positive prediction probability} - independent of correctness!
	\end{frame}
	
	\begin{frame}{Q27: Fairness Measures Summary Table 2 - Tutorial}
		\fontsize{6.5}{7.5}\selectfont
		\textbf{QUESTION: What question does equal opportunity answer?}
		
		\vspace{0.1cm}
		\textbf{OPTIONS:}
		\begin{enumerate}
			\item "Is the likelihood of a positive prediction the same across groups?"
			\item "Is the likelihood of a true positive prediction the same across groups?"
			\item "Is the likelihood of identifying actual positives as positive the same across groups?"
			\item "Is the likelihood of a negative prediction the same across groups?"
		\end{enumerate}
		
		\vspace{0.1cm}
		\textcolor{myblue}{\textbf{ANSWER: C}}
		
		\vspace{0.1cm}
		\textbf{DETAILED EXPLANATION:}
		\begin{itemize}
			\item \textbf{Why C is correct:} Equal opportunity asks whether the true positive rate (sensitivity) is the same across groups - whether actual positives are equally likely to be correctly identified regardless of group membership.
			\item \textbf{Why A is wrong:} This is demographic parity
			\item \textbf{Why B is wrong:} This is predictive rate parity
			\item \textbf{Why D is wrong:} This is not a standard fairness measure
		\end{itemize}
		
		\textcolor{myred}{\textbf{EXAM TRAP:}} Equal opportunity = \textcolor{myblue}{equal true positive rate} - "of actual positives, how many did we catch?"
	\end{frame}
	
	\begin{frame}{Q28: Base Rate Definition - Tutorial}
		\fontsize{6.5}{7.5}\selectfont
		\textbf{QUESTION: What is a base rate?}
		
		\vspace{0.1cm}
		\textbf{OPTIONS:}
		\begin{enumerate}
			\item The error rate of an algorithm
			\item The general prevalence of something within a population
			\item The accuracy of predictions
			\item The precision of predictions
		\end{enumerate}
		
		\vspace{0.1cm}
		\textcolor{myblue}{\textbf{ANSWER: B}}
		
		\vspace{0.1cm}
		\textbf{DETAILED EXPLANATION:}
		\begin{itemize}
			\item \textbf{Why B is correct:} Base rate refers to the general prevalence of something within a population. For example, the prevalence of a disease, or the proportion of defaulters. Base rates can differ significantly across groups.
			\item \textbf{Why A is wrong:} Error rate is different from base rate
			\item \textbf{Why C is wrong:} Accuracy is different from base rate
			\item \textbf{Why D is wrong:} Precision is different from base rate
		\end{itemize}
		
		\textcolor{myred}{\textbf{EXAM TRAP:}} Base rate = \textcolor{myblue}{prevalence in a population} - key for understanding fairness trade-offs!
	\end{frame}
	
	\begin{frame}{Q29: Performance vs Fairness Trade-off - Tutorial}
		\fontsize{6.5}{7.5}\selectfont
		\textbf{QUESTION: Why is balancing performance and fairness challenging?}
		
		\vspace{0.1cm}
		\textbf{OPTIONS:}
		\begin{enumerate}
			\item Fairness always improves performance
			\item Optimizing for fairness can degrade performance and vice versa
			\item They are always perfectly aligned
			\item Performance is never affected by fairness
		\end{enumerate}
		
		\vspace{0.1cm}
		\textcolor{myblue}{\textbf{ANSWER: B}}
		
		\vspace{0.1cm}
		\textbf{DETAILED EXPLANATION:}
		\begin{itemize}
			\item \textbf{Why B is correct:} There is often a fundamental trade-off between performance and fairness. Fairness constraints may reduce overall accuracy, and optimizing for performance may create unfair outcomes.
			\item \textbf{Why A is wrong:} Fairness often TRADES OFF with performance
			\item \textbf{Why C is wrong:} They are often in tension, not aligned
			\item \textbf{Why D is wrong:} Performance IS affected by fairness constraints
		\end{itemize}
		
		\textcolor{myred}{\textbf{EXAM TRAP:}} There is often a \textcolor{myblue}{trade-off between performance and fairness} - must choose based on context!
	\end{frame}
	
	\begin{frame}{Q30: Fairness Measures Summary - Tutorial}
		\fontsize{6.5}{7.5}\selectfont
		\textbf{QUESTION: Which fairness measure is most appropriate for medical diagnostics?}
		
		\vspace{0.1cm}
		\textbf{OPTIONS:}
		\begin{enumerate}
			\item Demographic parity
			\item Predictive rate parity
			\item Equal opportunity
			\item Accuracy parity
		\end{enumerate}
		
		\vspace{0.1cm}
		\textcolor{myblue}{\textbf{ANSWER: C}}
		
		\vspace{0.1cm}
		\textbf{DETAILED EXPLANATION:}
		\begin{itemize}
			\item \textbf{Why C is correct:} Equal opportunity ensures that positive cases (patients with disease) are equally likely to be caught across groups. Missing a disease (false negative) can be life-threatening, so equal sensitivity is critical.
			\item \textbf{Why A is wrong:} Demographic parity ignores prediction accuracy - not appropriate for medical diagnostics
			\item \textbf{Why B is wrong:} Predictive rate parity focuses on positive predictions, not on catching disease
			\item \textbf{Why D is wrong:} Accuracy may be misleading with imbalanced data
		\end{itemize}
		
		\textcolor{myred}{\textbf{EXAM TRAP:}} Medical diagnostics prioritizes \textcolor{myblue}{equal opportunity} (don't miss disease) - context determines the right measure!
	\end{frame}
	
	% ============================================================
	% SECTION 3: SOURCES OF UNFAIRNESS - QUESTIONS 31-50
	% ============================================================
	
	\begin{frame}{Q31: Problem Specification Bias - Tutorial}
		\fontsize{6.5}{7.5}\selectfont
		\textbf{QUESTION: Why can problem specification introduce bias?}
		
		\vspace{0.1cm}
		\textbf{OPTIONS:}
		\begin{enumerate}
			\item It always produces perfect algorithms
			\item How a problem is specified affects algorithmic performance and fairness
			\item Problem specification never affects outcomes
			\item All problems are specified identically
		\end{enumerate}
		
		\vspace{0.1cm}
		\textcolor{myblue}{\textbf{ANSWER: B}}
		
		\vspace{0.1cm}
		\textbf{DETAILED EXPLANATION:}
		\begin{itemize}
			\item \textbf{Why B is correct:} Choice of criteria affects who applies and who gets hired. Translation to measurable features requires proxies. Some concepts are harder to quantify than others (e.g., productivity vs. teamwork).
			\item \textbf{Why A is wrong:} Problem specification can introduce significant bias
			\item \textbf{Why C is wrong:} Problem specification significantly affects outcomes
			\item \textbf{Why D is wrong:} Problems vary widely in how they are specified
		\end{itemize}
		
		\textcolor{myred}{\textbf{EXAM TRAP:}} Problem specification is \textcolor{myblue}{critical for fairness} - how you define the problem matters!
	\end{frame}
	
	\begin{frame}{Q32: Proxy Variables - Tutorial}
		\fontsize{6.5}{7.5}\selectfont
		\textbf{QUESTION: What is a proxy variable in algorithmic design?}
		
		\vspace{0.1cm}
		\textbf{OPTIONS:}
		\begin{enumerate}
			\item A variable that perfectly measures the desired concept
			\item A variable used as a stand-in for a concept that is hard to measure directly
			\item A variable that is never used in algorithms
			\item A variable with no correlation to the desired concept
		\end{enumerate}
		
		\vspace{0.1cm}
		\textcolor{myblue}{\textbf{ANSWER: B}}
		
		\vspace{0.1cm}
		\textbf{DETAILED EXPLANATION:}
		\begin{itemize}
			\item \textbf{Why B is correct:} Proxies are used when direct measurement is difficult (e.g., productivity measured by sales). They reduce complex concepts to simpler measurable constructs but risk introducing bias if imperfect.
			\item \textbf{Why A is wrong:} Proxies are imperfect by definition
			\item \textbf{Why C is wrong:} Proxies ARE used in algorithms
			\item \textbf{Why D is wrong:} Proxies should be correlated with the concept
		\end{itemize}
		
		\textcolor{myred}{\textbf{EXAM TRAP:}} Proxies are \textcolor{myblue}{imperfect stand-ins} for hard-to-measure concepts - can introduce bias!
	\end{frame}
	
	\begin{frame}{Q33: Amazon Recruiting Tool Example - Tutorial}
		\fontsize{6.5}{7.5}\selectfont
		\textbf{QUESTION: What happened with Amazon's recruiting tool?}
		
		\vspace{0.1cm}
		\textbf{OPTIONS:}
		\begin{enumerate}
			\item It worked perfectly and was deployed
			\item It systematically downgraded CVs from women applicants
			\item It eliminated all bias
			\item It was immediately successful
		\end{enumerate}
		
		\vspace{0.1cm}
		\textcolor{myblue}{\textbf{ANSWER: B}}
		
		\vspace{0.1cm}
		\textbf{DETAILED EXPLANATION:}
		\begin{itemize}
			\item \textbf{Why B is correct:} Amazon attempted to create an unbiased recruiting tool but the algorithm systematically downgraded women's CVs. Even identical CVs were rated lower for women. Amazon eventually abandoned the project.
			\item \textbf{Why A is wrong:} The tool failed due to bias
			\item \textbf{Why C is wrong:} It perpetuated existing biases
			\item \textbf{Why D is wrong:} The project was abandoned
		\end{itemize}
		
		\textcolor{myred}{\textbf{EXAM TRAP:}} Amazon's tool \textcolor{myblue}{systematically discriminated against women} - demonstrates difficulty of fair AI!
	\end{frame}
	
	\begin{frame}{Q34: Fairness Through Unawareness - Tutorial}
		\fontsize{6.5}{7.5}\selectfont
		\textbf{QUESTION: What is "fairness through unawareness"?}
		
		\vspace{0.1cm}
		\textbf{OPTIONS:}
		\begin{enumerate}
			\item Using all available features
			\item Removing protected characteristics from input features
			\item Always including protected characteristics
			\item Ignoring all features
		\end{enumerate}
		
		\vspace{0.1cm}
		\textcolor{myblue}{\textbf{ANSWER: B}}
		
		\vspace{0.1cm}
		\textbf{DETAILED EXPLANATION:}
		\begin{itemize}
			\item \textbf{Why B is correct:} Fairness through unawareness is the intuitive approach of removing protected characteristics from input features, assuming this prevents discrimination. However, correlated variables can still carry protected information.
			\item \textbf{Why A is wrong:} This would increase bias risk
			\item \textbf{Why C is wrong:} This would allow explicit discrimination
			\item \textbf{Why D is wrong:} This would make predictions impossible
		\end{itemize}
		
		\textcolor{myred}{\textbf{EXAM TRAP:}} Fairness through unawareness \textcolor{myblue}{does NOT eliminate bias} - correlated variables still carry protected information!
	\end{frame}
	
	\begin{frame}{Q35: FTA Prediction Example - Tutorial}
		\fontsize{6.5}{7.5}\selectfont
		\textbf{QUESTION: In the FTA prediction example, what happened when protected characteristics were removed?}
		
		\vspace{0.1cm}
		\textbf{OPTIONS:}
		\begin{enumerate}
			\item Bias was completely eliminated
			\item The coefficient on prior convictions increased from 54% to 72%
			\item The algorithm became perfectly fair
			\item No changes occurred
		\end{enumerate}
		
		\vspace{0.1cm}
		\textcolor{myblue}{\textbf{ANSWER: B}}
		
		\vspace{0.1cm}
		\textbf{DETAILED EXPLANATION:}
		\begin{itemize}
			\item \textbf{Why B is correct:} The protected variable and prior convictions were correlated. Removing the protected variable contaminated the prior variable, increasing its coefficient from 54% to 72%. The algorithm still indirectly discriminated.
			\item \textbf{Why A is wrong:} Bias was NOT eliminated
			\item \textbf{Why C is wrong:} The algorithm remained biased
			\item \textbf{Why D is wrong:} Significant changes occurred
		\end{itemize}
		
		\textcolor{myred}{\textbf{EXAM TRAP:}} Removing protected characteristics \textcolor{myblue}{can make bias worse} through correlated variables!
	\end{frame}
	
	\begin{frame}{Q36: Historical Bias - Tutorial}
		\fontsize{6.5}{7.5}\selectfont
		\textbf{QUESTION: What is historical bias in AI?}
		
		\vspace{0.1cm}
		\textbf{OPTIONS:}
		\begin{enumerate}
			\item Bias that occurs only in the present
			\item Bias in training data that reflects past discrimination
			\item Bias that never affects algorithms
			\item Bias that is always intentional
		\end{enumerate}
		
		\vspace{0.1cm}
		\textcolor{myblue}{\textbf{ANSWER: B}}
		
		\vspace{0.1cm}
		\textbf{DETAILED EXPLANATION:}
		\begin{itemize}
			\item \textbf{Why B is correct:} Supervised learning algorithms learn from historical data, which reflects past biases and discrimination. If women were historically less likely selected for certain roles, the algorithm learns and may amplify this pattern.
			\item \textbf{Why A is wrong:} Historical bias is about the PAST
			\item \textbf{Why C is wrong:} Historical bias DOES affect algorithms
			\item \textbf{Why D is wrong:} Historical bias is often unintentional
		\end{itemize}
		
		\textcolor{myred}{\textbf{EXAM TRAP:}} Historical bias = \textcolor{myblue}{past discrimination reflected in training data} - algorithms learn from the past!
	\end{frame}
	
	\begin{frame}{Q37: Feedback Loops - Tutorial}
		\fontsize{6.5}{7.5}\selectfont
		\textbf{QUESTION: What is a feedback loop in AI systems?}
		
		\vspace{0.1cm}
		\textbf{OPTIONS:}
		\begin{enumerate}
			\item A system that corrects its own errors
			\item When algorithmic predictions influence future data, perpetuating bias
			\item A system that never changes
			\item A system that eliminates all bias
		\end{enumerate}
		
		\vspace{0.1cm}
		\textcolor{myblue}{\textbf{ANSWER: B}}
		
		\vspace{0.1cm}
		\textbf{DETAILED EXPLANATION:}
		\begin{itemize}
			\item \textbf{Why B is correct:} Feedback loops occur when algorithmic predictions influence real-world decisions, which create new data that reflects those biased decisions. The algorithm learns from this biased data, amplifying the original bias.
			\item \textbf{Why A is wrong:} Feedback loops often AMPLIFY errors
			\item \textbf{Why C is wrong:} Feedback loops cause change (often negative)
			\item \textbf{Why D is wrong:} Feedback loops can perpetuate bias
		\end{itemize}
		
		\textcolor{myred}{\textbf{EXAM TRAP:}} Feedback loops \textcolor{myblue}{amplify and perpetuate bias} - become self-reinforcing!
	\end{frame}
	
	\begin{frame}{Q38: Representation Bias - Tutorial}
		\fontsize{6.5}{7.5}\selectfont
		\textbf{QUESTION: What is representation bias?}
		
		\vspace{0.1cm}
		\textbf{OPTIONS:}
		\begin{enumerate}
			\item When data is perfectly representative
			\item When the collected sample is not representative of the target population
			\item When all groups are equally represented
			\item When data collection is unbiased
		\end{enumerate}
		
		\vspace{0.1cm}
		\textcolor{myblue}{\textbf{ANSWER: B}}
		
		\vspace{0.1cm}
		\textbf{DETAILED EXPLANATION:}
		\begin{itemize}
			\item \textbf{Why B is correct:} Representation bias occurs when the collected sample does not reflect the target population. Some groups are over-represented, others under-represented, leading to problematic model specifications.
			\item \textbf{Why A is wrong:} Representation bias means NOT representative
			\item \textbf{Why C is wrong:} This would be balanced representation
			\item \textbf{Why D is wrong:} Representation bias IS a form of bias
		\end{itemize}
		
		\textcolor{myred}{\textbf{EXAM TRAP:}} Representation bias = \textcolor{myblue}{non-representative sample} - algorithms perform better on dominant groups!
	\end{frame}
	
	\begin{frame}{Q39: Measurement Bias - Tutorial}
		\fontsize{6.5}{7.5}\selectfont
		\textbf{QUESTION: What is measurement bias?}
		
		\vspace{0.1cm}
		\textbf{OPTIONS:}
		\begin{enumerate}
			\item When measurement methods are consistent across the sample
			\item When measurement methods are not consistent across the sample
			\item When all measurements are perfectly accurate
			\item When measurements are always reliable
		\end{enumerate}
		
		\vspace{0.1cm}
		\textcolor{myblue}{\textbf{ANSWER: B}}
		
		\vspace{0.1cm}
		\textbf{DETAILED EXPLANATION:}
		\begin{itemize}
			\item \textbf{Why B is correct:} Measurement bias occurs when measurement methods are not consistent across the sample. Different groups are measured differently, which can favor certain groups and create biased data.
			\item \textbf{Why A is wrong:} Measurement bias means INCONSISTENT methods
			\item \textbf{Why C is wrong:} Measurement bias means NOT perfectly accurate
			\item \textbf{Why D is wrong:} Measurement bias means UNRELIABLE across groups
		\end{itemize}
		
		\textcolor{myred}{\textbf{EXAM TRAP:}} Measurement bias = \textcolor{myblue}{inconsistent measurement across groups} - creates systematic differences!
	\end{frame}
	
	\begin{frame}{Q40: Data Composition Bias - Tutorial}
		\fontsize{6.5}{7.5}\selectfont
		\textbf{QUESTION: What is data composition bias?}
		
		\vspace{0.1cm}
		\textbf{OPTIONS:}
		\begin{enumerate}
			\item When datasets are perfectly balanced
			\item When minority subgroups are underrepresented, causing algorithms to underperform
			\item When all groups are equally represented
			\item When data is always complete
		\end{enumerate}
		
		\vspace{0.1cm}
		\textcolor{myblue}{\textbf{ANSWER: B}}
		
		\vspace{0.1cm}
		\textbf{DETAILED EXPLANATION:}
		\begin{itemize}
			\item \textbf{Why B is correct:} Data composition bias occurs when minority subgroups are underrepresented in the data. The algorithm underperforms for these groups because there are fewer data points to learn from.
			\item \textbf{Why A is wrong:} Composition bias means IMBALANCED data
			\item \textbf{Why C is wrong:} This would be balanced data
			\item \textbf{Why D is wrong:} Composition bias means INCOMPLETE representation
		\end{itemize}
		
		\textcolor{myred}{\textbf{EXAM TRAP:}} Data composition bias = \textcolor{myblue}{underrepresented minority subgroups} - algorithms underperform for these groups!
	\end{frame}
	
	\begin{frame}{Q41: Balancing Datasets - Tutorial}
		\fontsize{6.5}{7.5}\selectfont
		\textbf{QUESTION: What does balancing a dataset mean?}
		
		\vspace{0.1cm}
		\textbf{OPTIONS:}
		\begin{enumerate}
			\item Using proportional representation
			\item Ensuring relevant minority subgroups occupy a similar proportion as majority groups
			\item Removing all minority data
			\item Always using the original distribution
		\end{enumerate}
		
		\vspace{0.1cm}
		\textcolor{myblue}{\textbf{ANSWER: B}}
		
		\vspace{0.1cm}
		\textbf{DETAILED EXPLANATION:}
		\begin{itemize}
			\item \textbf{Why B is correct:} Balancing means increasing the representation of minority groups to ensure the algorithm recognizes them as statistically relevant. This is different from proportional representation.
			\item \textbf{Why A is wrong:} Proportional representation is different from balancing
			\item \textbf{Why C is wrong:} Balancing INCREASES minority representation
			\item \textbf{Why D is wrong:} Balancing changes the original distribution
		\end{itemize}
		
		\textcolor{myred}{\textbf{EXAM TRAP:}} Balancing = \textcolor{myblue}{over-representing minorities} for algorithm performance - not proportional representation!
	\end{frame}
	
	\begin{frame}{Q42: Bootstrapping Methods - Tutorial}
		\fontsize{6.5}{7.5}\selectfont
		\textbf{QUESTION: What are methods for balancing datasets?}
		
		\vspace{0.1cm}
		\textbf{OPTIONS:}
		\begin{enumerate}
			\item Over-sampling minority groups
			\item Under-sampling majority groups
			\item Both A and B
			\item Only removing outliers
		\end{enumerate}
		
		\vspace{0.1cm}
		\textcolor{myblue}{\textbf{ANSWER: C}}
		
		\vspace{0.1cm}
		\textbf{DETAILED EXPLANATION:}
		\begin{itemize}
			\item \textbf{Why C is correct:} Both over-sampling (adding more data from minority groups) and under-sampling (removing data from majority groups) are methods for balancing datasets. They can be used separately or together.
			\item \textbf{Why A is wrong:} Correct but incomplete
			\item \textbf{Why B is wrong:} Correct but incomplete
			\item \textbf{Why D is wrong:} Outlier removal is a different technique
		\end{itemize}
		
		\textcolor{myred}{\textbf{EXAM TRAP:}} Bootstrapping methods include \textcolor{myblue}{over-sampling AND under-sampling} - often used together!
	\end{frame}
	
	\begin{frame}{Q43: Intersectional Bias - Tutorial}
		\fontsize{6.5}{7.5}\selectfont
		\textbf{QUESTION: What is intersectional bias?}
		
		\vspace{0.1cm}
		\textbf{OPTIONS:}
		\begin{enumerate}
			\item Bias that only affects one group
			\item Bias affecting groups at the intersection of multiple underrepresented subgroups
			\item Bias that never occurs
			\item Bias that is always intentional
		\end{enumerate}
		
		\vspace{0.1cm}
		\textcolor{myblue}{\textbf{ANSWER: B}}
		
		\vspace{0.1cm}
		\textbf{DETAILED EXPLANATION:}
		\begin{itemize}
			\item \textbf{Why B is correct:} Intersectional bias affects groups at the intersection of multiple underrepresented subgroups (e.g., women of color in facial recognition). Performance is particularly low for these intersectional groups.
			\item \textbf{Why A is wrong:} Intersectional bias affects MULTIPLE groups
			\item \textbf{Why C is wrong:} Intersectional bias DOES occur
			\item \textbf{Why D is wrong:} Intersectional bias is often unintentional
		\end{itemize}
		
		\textcolor{myred}{\textbf{EXAM TRAP:}} Intersectional bias = \textcolor{myblue}{multiple underrepresented dimensions combined} - hardest to detect and correct!
	\end{frame}
	
	\begin{frame}{Q44: Word Embeddings Bias - Tutorial}
		\fontsize{6.5}{7.5}\selectfont
		\textbf{QUESTION: How can word embeddings introduce bias?}
		
		\vspace{0.1cm}
		\textbf{OPTIONS:}
		\begin{enumerate}
			\item They always produce unbiased results
			\item They can contain stereotypes (e.g., man:computer programmer as woman:homemaker)
			\item They never contain stereotypes
			\item They are always perfectly neutral
		\end{enumerate}
		
		\vspace{0.1cm}
		\textcolor{myblue}{\textbf{ANSWER: B}}
		
		\vspace{0.1cm}
		\textbf{DETAILED EXPLANATION:}
		\begin{itemize}
			\item \textbf{Why B is correct:} Word embeddings capture meaning from text, which contains societal stereotypes. Research has shown that embeddings can reinforce stereotypes (e.g., "man is to computer programmer as woman is to homemaker").
			\item \textbf{Why A is wrong:} Word embeddings CAN contain bias
			\item \textbf{Why C is wrong:} Word embeddings DO contain stereotypes
			\item \textbf{Why D is wrong:} Word embeddings reflect societal biases
		\end{itemize}
		
		\textcolor{myred}{\textbf{EXAM TRAP:}} Word embeddings can contain \textcolor{myblue}{harmful stereotypes} - reflect societal biases in text!
	\end{frame}
	
	\begin{frame}{Q45: Objective Function Bias - Tutorial}
		\fontsize{6.5}{7.5}\selectfont
		\textbf{QUESTION: How can objective functions introduce bias?}
		
		\vspace{0.1cm}
		\textbf{OPTIONS:}
		\begin{enumerate}
			\item They never affect fairness
			\item What the algorithm optimizes for makes an ethically significant difference
			\item All objective functions are equally fair
			\item Objective functions always eliminate bias
		\end{enumerate}
		
		\vspace{0.1cm}
		\textcolor{myblue}{\textbf{ANSWER: B}}
		
		\vspace{0.1cm}
		\textbf{DETAILED EXPLANATION:}
		\begin{itemize}
			\item \textbf{Why B is correct:} The choice of objective function (what the algorithm optimizes for) has ethical significance. Different objectives (performance vs. error distribution) have different fairness implications depending on context.
			\item \textbf{Why A is wrong:} Objective functions significantly affect fairness
			\item \textbf{Why C is wrong:} Different objectives have different fairness implications
			\item \textbf{Why D is wrong:} Objective functions CAN introduce bias
		\end{itemize}
		
		\textcolor{myred}{\textbf{EXAM TRAP:}} Objective function choice = \textcolor{myblue}{ethical significance} - context determines appropriate objective!
	\end{frame}
	
	\begin{frame}{Q46: Cancer Detection Example - Tutorial}
		\fontsize{6.5}{7.5}\selectfont
		\textbf{QUESTION: In the cancer detection example, why might different algorithms require different approaches?}
		
		\vspace{0.1cm}
		\textbf{OPTIONS:}
		\begin{enumerate}
			\item All cancers are the same
			\item Skin cancer and brain cancer have different treatment risks and consequences
			\item Treatment is never invasive
			\item All patients are identical
		\end{enumerate}
		
		\vspace{0.1cm}
		\textcolor{myblue}{\textbf{ANSWER: B}}
		
		\vspace{0.1cm}
		\textbf{DETAILED EXPLANATION:}
		\begin{itemize}
			\item \textbf{Why B is correct:} Skin cancer (biopsy, relatively low risk) and brain cancer (major surgery, high risk) have different trade-offs for false positives/negatives. Context determines the appropriate objective function.
			\item \textbf{Why A is wrong:} Different cancers have different risks
			\item \textbf{Why C is wrong:} Some treatments ARE highly invasive
			\item \textbf{Why D is wrong:} Patients have different conditions
		\end{itemize}
		
		\textcolor{myred}{\textbf{EXAM TRAP:}} Context determines \textcolor{myblue}{appropriate trade-offs} - no one-size-fits-all objective function!
	\end{frame}
	
	\begin{frame}{Q47: Benchmark Dataset Bias - Tutorial}
		\fontsize{6.5}{7.5}\selectfont
		\textbf{QUESTION: How can benchmark datasets introduce bias?}
		
		\vspace{0.1cm}
		\textbf{OPTIONS:}
		\begin{enumerate}
			\item They are always representative
			\item They may not be representative of the intended use population
			\item They always eliminate bias
			\item They never affect model performance
		\end{enumerate}
		
		\vspace{0.1cm}
		\textcolor{myblue}{\textbf{ANSWER: B}}
		
		\vspace{0.1cm}
		\textbf{DETAILED EXPLANATION:}
		\begin{itemize}
			\item \textbf{Why B is correct:} Public benchmarks may not reflect real-world populations. Good performance on a benchmark doesn't guarantee good real-world performance and can mask biases that appear in deployment.
			\item \textbf{Why A is wrong:} Benchmarks may NOT be representative
			\item \textbf{Why C is wrong:} Benchmarks can CONTAIN bias
			\item \textbf{Why D is wrong:} Benchmarks significantly affect perceived performance
		\end{itemize}
		
		\textcolor{myred}{\textbf{EXAM TRAP:}} Benchmark datasets may \textcolor{myblue}{not reflect real-world populations} - test in the intended context!
	\end{frame}
	
	\begin{frame}{Q48: Deployment Context Bias - Tutorial}
		\fontsize{6.5}{7.5}\selectfont
		\textbf{QUESTION: How can deployment context introduce bias?}
		
		\vspace{0.1cm}
		\textbf{OPTIONS:}
		\begin{enumerate}
			\item Algorithms always perform equally in all contexts
			\item Algorithms used outside their designed context can lead to performance loss
			\item Context never matters
			\item All contexts are identical
		\end{enumerate}
		
		\vspace{0.1cm}
		\textcolor{myblue}{\textbf{ANSWER: B}}
		
		\vspace{0.1cm}
		\textbf{DETAILED EXPLANATION:}
		\begin{itemize}
			\item \textbf{Why B is correct:} Algorithms designed for one context may not work in another (e.g., recidivism prediction for sentence length). Historical data may be outdated, and user values may differ from developers.
			\item \textbf{Why A is wrong:} Performance varies by context
			\item \textbf{Why C is wrong:} Context significantly matters
			\item \textbf{Why D is wrong:} Contexts vary widely
		\end{itemize}
		
		\textcolor{myred}{\textbf{EXAM TRAP:}} Deployment context = \textcolor{myblue}{critical for performance and fairness} - validate in intended context!
	\end{frame}
	
	\begin{frame}{Q49: Human Oversight Importance - Tutorial}
		\fontsize{6.5}{7.5}\selectfont
		\textbf{QUESTION: Why is human oversight important for AI systems?}
		
		\vspace{0.1cm}
		\textbf{OPTIONS:}
		\begin{enumerate}
			\item Humans are never needed
			\item Ensuring safe and fair algorithms requires human oversight throughout
			\item AI systems are always perfectly safe
			\item Oversight is optional
		\end{enumerate}
		
		\vspace{0.1cm}
		\textcolor{myblue}{\textbf{ANSWER: B}}
		
		\vspace{0.1cm}
		\textbf{DETAILED EXPLANATION:}
		\begin{itemize}
			\item \textbf{Why B is correct:} Human oversight is essential throughout the AI lifecycle - data collection, model development, deployment, and regular auditing. This is critical in consequential domains.
			\item \textbf{Why A is wrong:} Human oversight is ESSENTIAL
			\item \textbf{Why C is wrong:} AI systems are NOT always perfectly safe
			\item \textbf{Why D is wrong:} Oversight is CRITICAL, not optional
		\end{itemize}
		
		\textcolor{myred}{\textbf{EXAM TRAP:}} Human oversight is \textcolor{myblue}{essential throughout the AI lifecycle} - not optional!
	\end{frame}
	
	\begin{frame}{Q50: Sources of Unfairness Summary - Tutorial}
		\fontsize{6.5}{7.5}\selectfont
		\textbf{QUESTION: Which is NOT a source of algorithmic unfairness?}
		
		\vspace{0.1cm}
		\textbf{OPTIONS:}
		\begin{enumerate}
			\item Problem specification and feature selection
			\item Data collection and composition
			\item Model development
			\item Always using perfect data
		\end{enumerate}
		
		\vspace{0.1cm}
		\textcolor{myblue}{\textbf{ANSWER: D}}
		
		\vspace{0.1cm}
		\textbf{DETAILED EXPLANATION:}
		\begin{itemize}
			\item \textbf{Why D is correct:} Perfect data does NOT exist - it's a hypothetical, not a source. All real data has limitations and potential biases.
			\item \textbf{Why A is wrong:} This IS a source of unfairness
			\item \textbf{Why B is wrong:} This IS a source of unfairness
			\item \textbf{Why C is wrong:} This IS a source of unfairness
		\end{itemize}
		
		\textcolor{myred}{\textbf{EXAM TRAP:}} Perfect data \textcolor{myblue}{does NOT exist} - all data has limitations that can introduce bias!
	\end{frame}
	
	% ============================================================
	% SECTION 4: EXPLAINABILITY, INTERPRETABILITY, TRANSPARENCY - QUESTIONS 51-65
	% ============================================================
	
	\begin{frame}{Q51: Explainability Definition - Tutorial}
		\fontsize{6.5}{7.5}\selectfont
		\textbf{QUESTION: What is explainability in AI?}
		
		\vspace{0.1cm}
		\textbf{OPTIONS:}
		\begin{enumerate}
			\item The ability to make predictions
			\item The capacity to explain in understandable terms how an AI system makes decisions
			\item The speed of the algorithm
			\item The accuracy of the algorithm
		\end{enumerate}
		
		\vspace{0.1cm}
		\textcolor{myblue}{\textbf{ANSWER: B}}
		
		\vspace{0.1cm}
		\textbf{DETAILED EXPLANATION:}
		\begin{itemize}
			\item \textbf{Why B is correct:} Explainability is the capacity to explain decisions in understandable terms, often after the fact (ex-post). It is essential for accountability and trust in high-stakes domains.
			\item \textbf{Why A is wrong:} This is prediction capability
			\item \textbf{Why C is wrong:} Speed is different from explainability
			\item \textbf{Why D is wrong:} Accuracy is different from explainability
		\end{itemize}
		
		\textcolor{myred}{\textbf{EXAM TRAP:}} Explainability = \textcolor{myblue}{explaining decisions after the fact} - essential for accountability!
	\end{frame}
	
	\begin{frame}{Q52: Interpretability Definition - Tutorial}
		\fontsize{6.5}{7.5}\selectfont
		\textbf{QUESTION: What is interpretability in AI?}
		
		\vspace{0.1cm}
		\textbf{OPTIONS:}
		\begin{enumerate}
			\item Explaining decisions after the fact
			\item The degree to which a human can comprehend and predict the model's behavior inherently
			\item The speed of the algorithm
			\item The accuracy of the algorithm
		\end{enumerate}
		
		\vspace{0.1cm}
		\textcolor{myblue}{\textbf{ANSWER: B}}
		
		\vspace{0.1cm}
		\textbf{DETAILED EXPLANATION:}
		\begin{itemize}
			\item \textbf{Why B is correct:} Interpretability is the degree to which a human can understand the model's behavior through built-in mechanisms (ex-ante). Inherently interpretable models like decision trees can be understood from the start.
			\item \textbf{Why A is wrong:} This is explainability
			\item \textbf{Why C is wrong:} Speed is different
			\item \textbf{Why D is wrong:} Accuracy is different
		\end{itemize}
		
		\textcolor{myred}{\textbf{EXAM TRAP:}} Interpretability = \textcolor{myblue}{inherently understandable models} - built-in, not after the fact!
	\end{frame}
	
	\begin{frame}{Q53: Transparency Definition - Tutorial}
		\fontsize{6.5}{7.5}\selectfont
		\textbf{QUESTION: What is transparency in AI?}
		
		\vspace{0.1cm}
		\textbf{OPTIONS:}
		\begin{enumerate}
			\item Making predictions quickly
			\item The openness and accessibility of an AI system's design, algorithms, data, and decision-making processes
			\item The accuracy of the algorithm
			\item The speed of the algorithm
		\end{enumerate}
		
		\vspace{0.1cm}
		\textcolor{myblue}{\textbf{ANSWER: B}}
		
		\vspace{0.1cm}
		\textbf{DETAILED EXPLANATION:}
		\begin{itemize}
			\item \textbf{Why B is correct:} Transparency is the openness and accessibility of an AI system's design, algorithms, data, and processes. It builds trust, enables oversight, and facilitates collaboration.
			\item \textbf{Why A is wrong:} Speed is different
			\item \textbf{Why C is wrong:} Accuracy is different
			\item \textbf{Why D is wrong:} Speed is different
		\end{itemize}
		
		\textcolor{myred}{\textbf{EXAM TRAP:}} Transparency = \textcolor{myblue}{openness and accessibility} - builds trust and enables oversight!
	\end{frame}
	
	\begin{frame}{Q54: Black-Box Problem - Tutorial}
		\fontsize{6.5}{7.5}\selectfont
		\textbf{QUESTION: What is the black-box problem in AI?}
		
		\vspace{0.1cm}
		\textbf{OPTIONS:}
		\begin{enumerate}
			\item Models that are always transparent
			\item Models that are too complex to be interpreted by humans
			\item Models that are always simple
			\item Models that are always accurate\end{enumerate}
		
		\vspace{0.1cm}
		\textcolor{myblue}{\textbf{ANSWER: B}}
		
		\vspace{0.1cm}
		\textbf{DETAILED EXPLANATION:}
		\begin{itemize}
			\item \textbf{Why B is correct:} The black-box problem occurs when models (especially neural networks) are too complex for humans to understand their intrinsic architecture. This creates explainability and accountability issues.
			\item \textbf{Why A is wrong:} Black boxes are NOT transparent
			\item \textbf{Why C is wrong:} Black boxes are COMPLEX
			\item \textbf{Why D is wrong:} Black boxes may NOT be accurate
		\end{itemize}
		
		\textcolor{myred}{\textbf{EXAM TRAP:}} Black-box problem = \textcolor{myblue}{models too complex to interpret} - creates accountability gaps!
	\end{frame}
	
	\begin{frame}{Q55: Feature Importance Scores - Tutorial}
		\fontsize{6.5}{7.5}\selectfont
		\textbf{QUESTION: What are feature importance scores used for?}
		
		\vspace{0.1cm}
		\textbf{OPTIONS:}
		\begin{enumerate}
			\item To make predictions faster
			\item To rank the importance of input features for outcomes
			\item To reduce the number of data points
			\item To increase model complexity
		\end{enumerate}
		
		\vspace{0.1cm}
		\textcolor{myblue}{\textbf{ANSWER: B}}
		
		\vspace{0.1cm}
		\textbf{DETAILED EXPLANATION:}
		\begin{itemize}
			\item \textbf{Why B is correct:} Feature importance scores rank the importance of input features for outcomes. Used in XAI, they reveal which factors have most influence and help identify potential biases.
			\item \textbf{Why A is wrong:} Feature importance is about interpretation
			\item \textbf{Why C is wrong:} It doesn't reduce data points
			\item \textbf{Why D is wrong:} It explains, not increases complexity
		\end{itemize}
		
		\textcolor{myred}{\textbf{EXAM TRAP:}} Feature importance = \textcolor{myblue}{ranks feature influence} - helps identify biases!
	\end{frame}
	
	\begin{frame}{Q56: Shapley Values - Tutorial}
		\fontsize{6.5}{7.5}\selectfont
		\textbf{QUESTION: What do Shapley values calculate?}
		
		\vspace{0.1cm}
		\textbf{OPTIONS:}
		\begin{enumerate}
			\item The accuracy of a model
			\item The contribution of each input feature to the model's output
			\item The speed of the model
			\item The size of the model
		\end{enumerate}
		
		\vspace{0.1cm}
		\textcolor{myblue}{\textbf{ANSWER: B}}
		
		\vspace{0.1cm}
		\textbf{DETAILED EXPLANATION:}
		\begin{itemize}
			\item \textbf{Why B is correct:} Shapley values calculate the contribution of each input feature by measuring the marginal effect of including each feature across all possible combinations. Based on game theory, they provide fair attribution.
			\item \textbf{Why A is wrong:} Accuracy is different
			\item \textbf{Why C is wrong:} Speed is different
			\item \textbf{Why D is wrong:} Size is different
		\end{itemize}
		
		\textcolor{myred}{\textbf{EXAM TRAP:}} Shapley values = \textcolor{myblue}{feature contribution calculation} - based on game theory!
	\end{frame}
	
	\begin{frame}{Q57: LUCID - Tutorial}
		\fontsize{6.5}{7.5}\selectfont
		\textbf{QUESTION: What does LUCID do in XAI?}
		
		\vspace{0.1cm}
		\textbf{OPTIONS:}
		\begin{enumerate}
			\item Makes predictions faster
			\item Works backward from algorithmic output to reveal internal logic
			\item Increases model complexity
			\item Removes features
		\end{enumerate}
		
		\vspace{0.1cm}
		\textcolor{myblue}{\textbf{ANSWER: B}}
		
		\vspace{0.1cm}
		\textbf{DETAILED EXPLANATION:}
		\begin{itemize}
			\item \textbf{Why B is correct:} LUCID (Locating Unfairness through Canonical Inverse Design) works backward from a particular output to create a distribution over the input space, revealing the model's internal logic and locating potential unfairness.
			\item \textbf{Why A is wrong:} Speed is not LUCID's purpose
			\item \textbf{Why C is wrong:} LUCID helps understand complexity
			\item \textbf{Why D is wrong:} LUCID doesn't remove features
		\end{itemize}
		
		\textcolor{myred}{\textbf{EXAM TRAP:}} LUCID works \textcolor{myblue}{backward from output to reveal logic} - locates unfairness!
	\end{frame}
	
	\begin{frame}{Q58: Surrogate Models - Tutorial}
		\fontsize{6.5}{7.5}\selectfont
		\textbf{QUESTION: What are surrogate models?}
		
		\vspace{0.1cm}
		\textbf{OPTIONS:}
		\begin{enumerate}
			\item More complex models that replace simpler ones
			\item Simpler, more interpretable models that approximate complex AI systems
			\item Models that are always less accurate
			\item Models that cannot be interpreted
		\end{enumerate}
		
		\vspace{0.1cm}
		\textcolor{myblue}{\textbf{ANSWER: B}}
		
		\vspace{0.1cm}
		\textbf{DETAILED EXPLANATION:}
		\begin{itemize}
			\item \textbf{Why B is correct:} Surrogate models are simpler, more interpretable models (like linear regression or decision trees) that approximate the behavior of complex AI systems to enhance interpretability.
			\item \textbf{Why A is wrong:} Surrogate models are SIMPLER
			\item \textbf{Why C is wrong:} Surrogate models can be quite accurate locally
			\item \textbf{Why D is wrong:} Surrogate models are INTERPRETABLE
		\end{itemize}
		
		\textcolor{myred}{\textbf{EXAM TRAP:}} Surrogate models are \textcolor{myblue}{simpler, interpretable approximations} - used for explanation!
	\end{frame}
	
	\begin{frame}{Q59: LIME Definition - Tutorial}
		\fontsize{6.5}{7.5}\selectfont
		\textbf{QUESTION: What is LIME in XAI?}
		
		\vspace{0.1cm}
		\textbf{OPTIONS:}
		\begin{enumerate}
			\item Local Interpretable Model-agnostic Explanations
			\item A method to make models faster
			\item A technique to increase accuracy
			\item A method to remove bias
		\end{enumerate}
		
		\vspace{0.1cm}
		\textcolor{myblue}{\textbf{ANSWER: A}}
		
		\vspace{0.1cm}
		\textbf{DETAILED EXPLANATION:}
		\begin{itemize}
			\item \textbf{Why A is correct:} LIME is Local Interpretable Model-agnostic Explanations. It approximates a complex model's predictions locally with an interpretable model, providing explanations that are specific to individual predictions.
			\item \textbf{Why B is wrong:} LIME is about explanation, not speed
			\item \textbf{Why C is wrong:} LIME is about interpretability, not accuracy
			\item \textbf{Why D is wrong:} LIME helps identify bias, doesn't remove it directly
		\end{itemize}
		
		\textcolor{myred}{\textbf{EXAM TRAP:}} LIME = \textcolor{myblue}{Local Interpretable Model-agnostic Explanations} - explains individual predictions!
	\end{frame}
	
	\begin{frame}{Q60: XAI Challenges - Tutorial}
		\fontsize{6.5}{7.5}\selectfont
		\textbf{QUESTION: What is a challenge with XAI techniques?}
		
		\vspace{0.1cm}
		\textbf{OPTIONS:}
		\begin{enumerate}
			\item They always produce perfect explanations
			\item Oversimplification can distort understanding
			\item They are never computationally expensive
			\item They always eliminate bias
		\end{enumerate}
		
		\vspace{0.1cm}
		\textcolor{myblue}{\textbf{ANSWER: B}}
		
		\vspace{0.1cm}
		\textbf{DETAILED EXPLANATION:}
		\begin{itemize}
			\item \textbf{Why B is correct:} Oversimplification of complex decisions can distort our understanding and lead to misinterpretation. Additionally, XAI techniques can be computationally expensive and are ex-post.
			\item \textbf{Why A is wrong:} Explanations can be imperfect
			\item \textbf{Why C is wrong:} XAI CAN be computationally expensive
			\item \textbf{Why D is wrong:} XAI helps identify, not eliminate bias
		\end{itemize}
		
		\textcolor{myred}{\textbf{EXAM TRAP:}} XAI can \textcolor{myblue}{oversimplify and distort understanding} - a significant challenge!
	\end{frame}
	
	\begin{frame}{Q61: Secrecy - Tutorial}
		\fontsize{6.5}{7.5}\selectfont
		\textbf{QUESTION: What is secrecy as a form of opaqueness?}
		
		\vspace{0.1cm}
		\textbf{OPTIONS:}
		\begin{enumerate}
			\item When algorithms are fully transparent
			\item When individuals are not aware of being subject to algorithmic decision making
			\item When algorithms are always explained
			\item When all data is public
		\end{enumerate}
		
		\vspace{0.1cm}
		\textcolor{myblue}{\textbf{ANSWER: B}}
		
		\vspace{0.1cm}
		\textbf{DETAILED EXPLANATION:}
		\begin{itemize}
			\item \textbf{Why B is correct:} Secrecy occurs when individuals are not aware of being subject to algorithmic profiling. Data is extracted without their knowledge (e.g., social media data used for targeted advertising).
			\item \textbf{Why A is wrong:} Secrecy is the OPPOSITE of transparent
			\item \textbf{Why C is wrong:} Secrecy means NOT explained
			\item \textbf{Why D is wrong:} Secrecy means NOT public
		\end{itemize}
		
		\textcolor{myred}{\textbf{EXAM TRAP:}} Secrecy = \textcolor{myblue}{unawareness of algorithmic decision making} - violates transparency and consent!
	\end{frame}
	
	\begin{frame}{Q62: Confidentiality - Tutorial}
		\fontsize{6.5}{7.5}\selectfont
		\textbf{QUESTION: What is confidentiality as a form of opaqueness?}
		
		\vspace{0.1cm}
		\textbf{OPTIONS:}
		\begin{enumerate}
			\item When algorithms are completely transparent
			\item When the algorithmic decision-making process is opaque to affected parties
			\item When all data is public
			\item When algorithms are always explained
		\end{enumerate}
		
		\vspace{0.1cm}
		\textcolor{myblue}{\textbf{ANSWER: B}}
		
		\vspace{0.1cm}
		\textbf{DETAILED EXPLANATION:}
		\begin{itemize}
			\item \textbf{Why B is correct:} Confidentiality occurs when affected parties know they're subject to AI but don't have knowledge of how decisions are made. This is a main concern in the AI discourse.
			\item \textbf{Why A is wrong:} Confidentiality means NOT transparent
			\item \textbf{Why C is wrong:} Confidentiality means NOT public
			\item \textbf{Why D is wrong:} Confidentiality means NOT explained
		\end{itemize}
		
		\textcolor{myred}{\textbf{EXAM TRAP:}} Confidentiality = \textcolor{myblue}{opaque decision-making process} - can't understand why decisions were made!
	\end{frame}
	
	\begin{frame}{Q63: Specialized Knowledge - Tutorial}
		\fontsize{6.5}{7.5}\selectfont
		\textbf{QUESTION: What is specialized knowledge as a form of opaqueness?}
		
		\vspace{0.1cm}
		\textbf{OPTIONS:}
		\begin{enumerate}
			\item When everyone can understand algorithms
			\item When understanding algorithms requires specialized training
			\item When algorithms are always simple
			\item When no knowledge is needed
		\end{enumerate}
		
		\vspace{0.1cm}
		\textcolor{myblue}{\textbf{ANSWER: B}}
		
		\vspace{0.1cm}
		\textbf{DETAILED EXPLANATION:}
		\begin{itemize}
			\item \textbf{Why B is correct:} Specialized knowledge opacity occurs when even with access to source code and data, understanding requires specialized training. This creates unequal access to understanding.
			\item \textbf{Why A is wrong:} Specialized knowledge means NOT everyone understands
			\item \textbf{Why C is wrong:} Algorithms may be complex
			\item \textbf{Why D is wrong:} Knowledge IS needed
		\end{itemize}
		
		\textcolor{myred}{\textbf{EXAM TRAP:}} Specialized knowledge = \textcolor{myblue}{requires training to understand} - not unique to AI!
	\end{frame}
	
	\begin{frame}{Q64: Reasons for Opacity - Tutorial}
		\fontsize{6.5}{7.5}\selectfont
		\textbf{QUESTION: What is a valid reason for algorithmic opacity?}
		
		\vspace{0.1cm}
		\textbf{OPTIONS:}
		\begin{enumerate}
			\item To hide discrimination
			\item Security (spam filters) and proprietary interests
			\item To avoid accountability
			\item To confuse users
		\end{enumerate}
		
		\vspace{0.1cm}
		\textcolor{myblue}{\textbf{ANSWER: B}}
		
		\vspace{0.1cm}
		\textbf{DETAILED EXPLANATION:}
		\begin{itemize}
			\item \textbf{Why B is correct:} Security (making code public enables circumvention) and proprietary interests (trade secrets) are legitimate reasons for some opacity. Must balance transparency with these interests.
			\item \textbf{Why A is wrong:} Hiding discrimination is NOT valid
			\item \textbf{Why C is wrong:} Avoiding accountability is NOT valid
			\item \textbf{Why D is wrong:} Confusing users is NOT valid
		\end{itemize}
		
		\textcolor{myred}{\textbf{EXAM TRAP:}} Legitimate reasons include \textcolor{myblue}{security and proprietary interests} - not hiding discrimination!
	\end{frame}
	
	\begin{frame}{Q65: Explainability Summary - Tutorial}
		\fontsize{6.5}{7.5}\selectfont
		\textbf{QUESTION: What is the most comprehensive statement about explainability?}
		
		\vspace{0.1cm}
		\textbf{OPTIONS:}
		\begin{enumerate}
			\item Explainability is always easy to achieve
			\item Explainability, interpretability, and transparency are essential for accountability and trust in AI systems
			\item Explainability is never needed
			\item Only black-box models are used
		\end{enumerate}
		
		\vspace{0.1cm}
		\textcolor{myblue}{\textbf{ANSWER: B}}
		
		\vspace{0.1cm}
		\textbf{DETAILED EXPLANATION:}
		\begin{itemize}
			\item \textbf{Why B is correct:} Explainability (ex-post understanding), interpretability (inherent understanding), and transparency (openness) are all essential for accountability and trust in AI systems, especially in high-stakes domains.
			\item \textbf{Why A is wrong:} Explainability can be challenging
			\item \textbf{Why C is wrong:} Explainability IS needed for accountability
			\item \textbf{Why D is wrong:} Many models are interpretable
		\end{itemize}
		
		\textcolor{myred}{\textbf{EXAM TRAP:}} All three are \textcolor{myblue}{essential for accountability and trust} - not optional!
	\end{frame}
	
	% ============================================================
	% SECTION 5: AUTONOMY AND MANIPULATION - QUESTIONS 66-70
	% ============================================================
	
	\begin{frame}{Q66: Human Autonomy Definition - Tutorial}
		\fontsize{6.5}{7.5}\selectfont
		\textbf{QUESTION: What is human autonomy in the context of AI?}
		
		\vspace{0.1cm}
		\textbf{OPTIONS:}
		\begin{enumerate}
			\item The ability to make any decision
			\item The ability to hold beliefs, make, and execute decisions of our own
			\item The ability to control AI systems
			\item The ability to avoid all technology
		\end{enumerate}
		
		\vspace{0.1cm}
		\textcolor{myblue}{\textbf{ANSWER: B}}
		
		\vspace{0.1cm}
		\textbf{DETAILED EXPLANATION:}
		\begin{itemize}
			\item \textbf{Why B is correct:} Human autonomy has two dimensions: internal (values, beliefs, decisions that are our own) and external (freedom and opportunity to execute decisions). It requires meaningful human oversight.
			\item \textbf{Why A is wrong:} Autonomy is more than making any decision
			\item \textbf{Why C is wrong:} This is control, not autonomy
			\item \textbf{Why D is wrong:} Avoiding technology is not autonomy
		\end{itemize}
		
		\textcolor{myred}{\textbf{EXAM TRAP:}} Autonomy = \textcolor{myblue}{own beliefs and decisions + ability to execute} - not just control!
	\end{frame}
	
	\begin{frame}{Q67: Cambridge Analytica Scandal - Tutorial}
		\fontsize{6.5}{7.5}\selectfont
		\textbf{QUESTION: What did the Cambridge Analytica scandal demonstrate?}
		
		\vspace{0.1cm}
		\textbf{OPTIONS:}
		\begin{enumerate}
			\item The benefits of data collection
			\item How big data could be used to manipulate voter behavior
			\item That data privacy is not important
			\item That AI never affects elections
		\end{enumerate}
		
		\vspace{0.1cm}
		\textcolor{myblue}{\textbf{ANSWER: B}}
		
		\vspace{0.1cm}
		\textbf{DETAILED EXPLANATION:}
		\begin{itemize}
			\item \textbf{Why B is correct:} Cambridge Analytica collected Facebook user data, inferred personality profiles from online behavior, and used this for psychologically tailored political advertising to manipulate voter behavior.
			\item \textbf{Why A is wrong:} The scandal showed MISUSE of data
			\item \textbf{Why C is wrong:} The scandal highlighted privacy concerns
			\item \textbf{Why D is wrong:} AI CAN affect elections
		\end{itemize}
		
		\textcolor{myred}{\textbf{EXAM TRAP:}} Cambridge Analytica = \textcolor{myblue}{political manipulation through data} - powerful demonstration of AI misuse!
	\end{frame}
	
	\begin{frame}{Q68: Facebook Emotional Contagion Experiment - Tutorial}
		\fontsize{6.5}{7.5}\selectfont
		\textbf{QUESTION: What did the Facebook emotional contagion experiment show?}
		
		\vspace{0.1cm}
		\textbf{OPTIONS:}
		\begin{enumerate}
			\item That emotions cannot be transmitted online
			\item That recommendation algorithms can shape and manipulate user experience
			\item That users always consent to experiments
			\item That AI has no effect on emotions
		\end{enumerate}
		
		\vspace{0.1cm}
		\textcolor{myblue}{\textbf{ANSWER: B}}
		
		\vspace{0.1cm}
		\textbf{DETAILED EXPLANATION:}
		\begin{itemize}
			\item \textbf{Why B is correct:} Researchers changed users' timelines to display predominantly positive or negative posts, and users' own posts became more positive/negative. This demonstrated the power of algorithms to shape user experience.
			\item \textbf{Why A is wrong:} Emotions CAN be transmitted online
			\item \textbf{Why C is wrong:} Users did NOT consent
			\item \textbf{Why D is wrong:} AI CAN affect emotions
		\end{itemize}
		
		\textcolor{myred}{\textbf{EXAM TRAP:}} Experiment showed \textcolor{myblue}{AI can manipulate emotional state} - without user consent!
	\end{frame}
	
	\begin{frame}{Q69: AI Manipulation Prevention - Tutorial}
		\fontsize{6.5}{7.5}\selectfont
		\textbf{QUESTION: How can AI manipulation be prevented?}
		
		\vspace{0.1cm}
		\textbf{OPTIONS:}
		\begin{enumerate}
			\item By ignoring all risks
			\item Through human oversight and regulatory oversight
			\item By eliminating all AI
			\item By never collecting data
		\end{enumerate}
		
		\vspace{0.1cm}
		\textcolor{myblue}{\textbf{ANSWER: B}}
		
		\vspace{0.1cm}
		\textbf{DETAILED EXPLANATION:}
		\begin{itemize}
			\item \textbf{Why B is correct:} Prevention requires human oversight throughout development and deployment, regulatory oversight (e.g., EU AI Act), transparency obligations, and balancing benefits with risks.
			\item \textbf{Why A is wrong:} Ignoring risks is not prevention
			\item \textbf{Why C is wrong:} Eliminating AI is unrealistic
			\item \textbf{Why D is wrong:} Some data collection is necessary
		\end{itemize}
		
		\textcolor{myred}{\textbf{EXAM TRAP:}} Prevention requires \textcolor{myblue}{human and regulatory oversight} - not ignoring or eliminating!
	\end{frame}
	
	\begin{frame}{Q70: Autonomy and Manipulation Summary - Tutorial}
		\fontsize{6.5}{7.5}\selectfont
		\textbf{QUESTION: What is the most comprehensive statement about autonomy and manipulation?}
		
		\vspace{0.1cm}
		\textbf{OPTIONS:}
		\begin{enumerate}
			\item AI never affects autonomy
			\item AI systems can shape our experience and potentially manipulate us, requiring oversight and regulation
			\item Manipulation is always beneficial
			\item Autonomy is not important
		\end{enumerate}
		
		\vspace{0.1cm}
		\textcolor{myblue}{\textbf{ANSWER: B}}
		
		\vspace{0.1cm}
		\textbf{DETAILED EXPLANATION:}
		\begin{itemize}
			\item \textbf{Why B is correct:} AI shapes online experience and can influence behavior through algorithms (Cambridge Analytica, emotional contagion). This requires oversight and regulation to prevent manipulation.
			\item \textbf{Why A is wrong:} AI CAN affect autonomy
			\item \textbf{Why C is wrong:} Manipulation is generally harmful
			\item \textbf{Why D is wrong:} Autonomy IS important
		\end{itemize}
		
		\textcolor{myred}{\textbf{EXAM TRAP:}} AI can \textcolor{myblue}{shape experience and manipulate} - requires robust oversight!
	\end{frame}
	
	% ============================================================
	% SECTION 6: SAFETY AND WELL-BEING - QUESTIONS 71-75
	% ============================================================
	
	\begin{frame}{Q71: Uber Crash Example - Tutorial}
		\fontsize{6.5}{7.5}\selectfont
		\textbf{QUESTION: What did the Uber driverless vehicle crash demonstrate?}
		
		\vspace{0.1cm}
		\textbf{OPTIONS:}
		\begin{enumerate}
			\item That autonomous vehicles are always safe
			\item The risks of over-dependence on AI in critical systems
			\item That safety mechanisms are never needed
			\item That AI never fails
		\end{enumerate}
		
		\vspace{0.1cm}
		\textcolor{myblue}{\textbf{ANSWER: B}}
		
		\vspace{0.1cm}
		\textbf{DETAILED EXPLANATION:}
		\begin{itemize}
			\item \textbf{Why B is correct:} The crash demonstrated multiple failures: vehicle systems failed to recognize the pedestrian, safety mechanisms were deactivated, and the human safety driver was distracted (automation bias). This shows risks of over-reliance on AI.
			\item \textbf{Why A is wrong:} The crash showed they are NOT always safe
			\item \textbf{Why C is wrong:} Safety mechanisms ARE needed
			\item \textbf{Why D is wrong:} AI CAN fail
		\end{itemize}
		
		\textcolor{myred}{\textbf{EXAM TRAP:}} Uber crash = \textcolor{myblue}{risks of over-reliance on AI} - demonstrates multiple failure points!
	\end{frame}
	
	\begin{frame}{Q72: Automation Bias Definition - Tutorial}
		\fontsize{6.5}{7.5}\selectfont
		\textbf{QUESTION: What is automation bias?}
		
		\vspace{0.1cm}
		\textbf{OPTIONS:}
		\begin{enumerate}
			\item The tendency to distrust automated systems
			\item The tendency of humans to over-rely on automated systems
			\item The tendency to always use manual methods
			\item The tendency to ignore technology
		\end{enumerate}
		
		\vspace{0.1cm}
		\textcolor{myblue}{\textbf{ANSWER: B}}
		
		\vspace{0.1cm}
		\textbf{DETAILED EXPLANATION:}
		\begin{itemize}
			\item \textbf{Why B is correct:} Automation bias is the tendency of humans to over-rely on automated systems, leading to complacency, reduced vigilance, and errors of oversight. Particularly problematic in critical situations.
			\item \textbf{Why A is wrong:} This is distrust, not over-reliance
			\item \textbf{Why C is wrong:} Automation bias is about using automation
			\item \textbf{Why D is wrong:} Automation bias is about over-relying
		\end{itemize}
		
		\textcolor{myred}{\textbf{EXAM TRAP:}} Automation bias = \textcolor{myblue}{over-reliance on automated systems} - leads to complacency!
	\end{frame}
	
	\begin{frame}{Q73: Skill Atrophy - Tutorial}
		\fontsize{6.5}{7.5}\selectfont
		\textbf{QUESTION: What is skill atrophy in the context of AI?}
		
		\vspace{0.1cm}
		\textbf{OPTIONS:}
		\begin{enumerate}
			\item Improving skills through AI use
			\item Decay in essential skills from prolonged reliance on automation
			\item Learning new skills quickly
			\item Maintaining all skills
		\end{enumerate}
		
		\vspace{0.1cm}
		\textcolor{myblue}{\textbf{ANSWER: B}}
		
		\vspace{0.1cm}
		\textbf{DETAILED EXPLANATION:}
		\begin{itemize}
			\item \textbf{Why B is correct:} Skill atrophy occurs when operators move from active decision makers to passive observers, causing essential skills to decay. This is a consequence of automation bias and prolonged automation use.
			\item \textbf{Why A is wrong:} This is skill improvement
			\item \textbf{Why C is wrong:} Skill atrophy is about DECAY, not learning
			\item \textbf{Why D is wrong:} Skill atrophy is about LOSING skills
		\end{itemize}
		
		\textcolor{myred}{\textbf{EXAM TRAP:}} Skill atrophy = \textcolor{myblue}{decay from over-reliance on automation} - loss of expertise!
	\end{frame}
	
	\begin{frame}{Q74: Automation Bias Mitigation - Tutorial}
		\fontsize{6.5}{7.5}\selectfont
		\textbf{QUESTION: How can automation bias be mitigated?}
		
		\vspace{0.1cm}
		\textbf{OPTIONS:}
		\begin{enumerate}
			\item By removing all automation
			\item By informing users and requiring periodic human input
			\item By ignoring the problem
			\item By increasing automation
		\end{enumerate}
		
		\vspace{0.1cm}
		\textcolor{myblue}{\textbf{ANSWER: B}}
		
		\vspace{0.1cm}
		\textbf{DETAILED EXPLANATION:}
		\begin{itemize}
			\item \textbf{Why B is correct:} Mitigation requires informing users about the risks and limitations of AI, and implementing design mechanisms that require periodic human input and verification (e.g., eye tracking with alerts).
			\item \textbf{Why A is wrong:} Removing all automation is unrealistic
			\item \textbf{Why C is wrong:} Ignoring the problem is not mitigation
			\item \textbf{Why D is wrong:} Increasing automation worsens the problem
		\end{itemize}
		
		\textcolor{myred}{\textbf{EXAM TRAP:}} Mitigation = \textcolor{myblue}{user education + periodic human verification} - keep operators engaged!
	\end{frame}
	
	\begin{frame}{Q75: Mental Health AI Risks - Tutorial}
		\fontsize{6.5}{7.5}\selectfont
		\textbf{QUESTION: What is a concern with AI-driven mental health support?}
		
		\vspace{0.1cm}
		\textbf{OPTIONS:}
		\begin{enumerate}
			\item AI systems have perfect empathy
			\item AI systems have no empathy, which is crucial for mental health support
			\item AI systems are always better than humans
			\item AI systems never fail in mental health
		\end{enumerate}
		
		\vspace{0.1cm}
		\textcolor{myblue}{\textbf{ANSWER: B}}
		
		\vspace{0.1cm}
		\textbf{DETAILED EXPLANATION:}
		\begin{itemize}
			\item \textbf{Why B is correct:} While AI can give seemingly empathetic responses, it has no genuine empathy - which is crucial for mental health support. Positive effects dwindle when patients know responses are AI-generated.
			\item \textbf{Why A is wrong:} AI does NOT have perfect empathy
			\item \textbf{Why C is wrong:} AI is NOT always better than humans
			\item \textbf{Why D is wrong:} AI CAN fail in mental health
		\end{itemize}
		
		\textcolor{myred}{\textbf{EXAM TRAP:}} AI lacks \textcolor{myblue}{genuine empathy} - critical for mental health support!
	\end{frame}
	
	% ============================================================
	% SECTION 7: REPUTATIONAL RISKS - QUESTIONS 76-80
	% ============================================================
	
	\begin{frame}{Q76: Reputational Risk Causes - Tutorial}
		\fontsize{6.5}{7.5}\selectfont
		\textbf{QUESTION: What are the three main causes of AI-related reputational damage?}
		
		\vspace{0.1cm}
		\textbf{OPTIONS:}
		\begin{enumerate}
			\item Speed, cost, and accuracy
			\item Privacy breaches, algorithmic bias, and lack of explainability
			\item Size, complexity, and speed
			\item Training data, model type, and deployment
		\end{enumerate}
		
		\vspace{0.1cm}
		\textcolor{myblue}{\textbf{ANSWER: B}}
		
		\vspace{0.1cm}
		\textbf{DETAILED EXPLANATION:}
		\begin{itemize}
			\item \textbf{Why B is correct:} Research shows the three main causes are privacy breaches (Cambridge Analytica), algorithmic bias (Dutch child benefits scandal), and lack of explainability (customers denied without reasons).
			\item \textbf{Why A is wrong:} Speed, cost, accuracy are not the main causes
			\item \textbf{Why C is wrong:} Size, complexity, speed are not the main causes
			\item \textbf{Why D is wrong:} Training data, model type, deployment are not the main causes
		\end{itemize}
		
		\textcolor{myred}{\textbf{EXAM TRAP:}} Three main causes = \textcolor{myblue}{privacy, bias, explainability} - critical to manage!
	\end{frame}
	
	\begin{frame}{Q77: Dutch Child Benefits Scandal - Tutorial}
		\fontsize{6.5}{7.5}\selectfont
		\textbf{QUESTION: What happened in the Dutch child benefits scandal?}
		
		\vspace{0.1cm}
		\textbf{OPTIONS:}
		\begin{enumerate}
			\item An algorithm correctly identified fraud
			\item A fraud-detection algorithm falsely denied benefits to thousands of citizens
			\item All claims were approved
			\item The algorithm worked perfectly
		\end{enumerate}
		
		\vspace{0.1cm}
		\textcolor{myblue}{\textbf{ANSWER: B}}
		
		\vspace{0.1cm}
		\textbf{DETAILED EXPLANATION:}
		\begin{itemize}
			\item \textbf{Why B is correct:} A fraud-detection algorithm deployed by the Dutch government falsely denied and reclaimed child benefits from thousands of citizens, with dire consequences. The scandal led to the Dutch prime minister's resignation.
			\item \textbf{Why A is wrong:} The algorithm incorrectly identified fraud
			\item \textbf{Why C is wrong:} Many claims were DENIED
			\item \textbf{Why D is wrong:} The algorithm failed catastrophically
		\end{itemize}
		
		\textcolor{myred}{\textbf{EXAM TRAP:}} Dutch scandal = \textcolor{myblue}{algorithmic bias with severe consequences} - prime minister resigned!
	\end{frame}
	
	\begin{frame}{Q78: Competence Criticism - Tutorial}
		\fontsize{6.5}{7.5}\selectfont
		\textbf{QUESTION: What is competence criticism of companies?}
		
		\vspace{0.1cm}
		\textbf{OPTIONS:}
		\begin{enumerate}
			\item Criticism that the company lacks relevant expertise
			\item Criticism that the company has misaligned values
			\item Criticism that the company is too successful
			\item Criticism that the company is too transparent
		\end{enumerate}
		
		\vspace{0.1cm}
		\textcolor{myblue}{\textbf{ANSWER: A}}
		
		\vspace{0.1cm}
		\textbf{DETAILED EXPLANATION:}
		\begin{itemize}
			\item \textbf{Why A is correct:} Competence criticism arises when there is a perception that the company lacks the relevant expertise and capacities to ensure their algorithms are safe, fair, and explainable.
			\item \textbf{Why B is wrong:} This is value alignment criticism
			\item \textbf{Why C is wrong:} This is not criticism
			\item \textbf{Why D is wrong:} Transparency is generally positive
		\end{itemize}
		
		\textcolor{myred}{\textbf{EXAM TRAP:}} Competence = \textcolor{myblue}{lack of expertise or capacity} - can be fixed through improvement!
	\end{frame}
	
	\begin{frame}{Q79: Value Alignment Criticism - Tutorial}
		\fontsize{6.5}{7.5}\selectfont
		\textbf{QUESTION: What is value alignment criticism of companies?}
		
		\vspace{0.1cm}
		\textbf{OPTIONS:}
		\begin{enumerate}
			\item Criticism that the company lacks relevant expertise
			\item Criticism that the company has misaligned values or acted despite knowing risks
			\item Criticism that the company is too successful
			\item Criticism that the company is too transparent
		\end{enumerate}
		
		\vspace{0.1cm}
		\textcolor{myblue}{\textbf{ANSWER: B}}
		
		\vspace{0.1cm}
		\textbf{DETAILED EXPLANATION:}
		\begin{itemize}
			\item \textbf{Why B is correct:} Value alignment criticism occurs when the company was fully aware of the risks but nevertheless proceeded (e.g., Facebook emotional contagion experiment). It suggests a violation of user trust.
			\item \textbf{Why A is wrong:} This is competence criticism
			\item \textbf{Why C is wrong:} This is not criticism
			\item \textbf{Why D is wrong:} Transparency is generally positive
		\end{itemize}
		
		\textcolor{myred}{\textbf{EXAM TRAP:}} Value alignment = \textcolor{myblue}{misaligned values or knowingly proceeding with risks} - harder to fix!
	\end{frame}
	
	\begin{frame}{Q80: Reputational Risk Mitigation - Tutorial}
		\fontsize{6.5}{7.5}\selectfont
		\textbf{QUESTION: What is NOT a reputational risk mitigation strategy?}
		
		\vspace{0.1cm}
		\textbf{OPTIONS:}
		\begin{enumerate}
			\item Continuous monitoring and evaluation
			\item Ignoring incidents and hoping they go away
			\item Ethical AI frameworks
			\item Stakeholder engagement and communication
		\end{enumerate}
		
		\vspace{0.1cm}
		\textcolor{myblue}{\textbf{ANSWER: B}}
		
		\vspace{0.1cm}
		\textbf{DETAILED EXPLANATION:}
		\begin{itemize}
			\item \textbf{Why B is correct:} Ignoring incidents is NOT a valid mitigation strategy. Effective strategies include continuous monitoring, ethical AI frameworks, stakeholder engagement, and crisis management plans.
			\item \textbf{Why A is wrong:} This IS a mitigation strategy
			\item \textbf{Why C is wrong:} This IS a mitigation strategy
			\item \textbf{Why D is wrong:} This IS a mitigation strategy
		\end{itemize}
		
		\textcolor{myred}{\textbf{EXAM TRAP:}} Ignoring incidents is \textcolor{myblue}{NOT a valid mitigation strategy} - proactive management is essential!
	\end{frame}
	
	% ============================================================
	% SECTION 8: EXISTENTIAL AND GLOBAL RISKS - QUESTIONS 81-85
	% ============================================================
	
	\begin{frame}{Q81: Superintelligence Definition - Tutorial}
		\fontsize{6.5}{7.5}\selectfont
		\textbf{QUESTION: What is superintelligence?}
		
		\vspace{0.1cm}
		\textbf{OPTIONS:}
		\begin{enumerate}
			\item An AI system that is faster than humans
			\item An AI system that operates beyond human-level intelligence
			\item An AI system that is less intelligent than humans
			\item An AI system that only performs one task
		\end{enumerate}
		
		\vspace{0.1cm}
		\textcolor{myblue}{\textbf{ANSWER: B}}
		
		\vspace{0.1cm}
		\textbf{DETAILED EXPLANATION:}
		\begin{itemize}
			\item \textbf{Why B is correct:} Superintelligence refers to AI that operates beyond human-level intelligence. Such a system could act in ways that threaten humanity, especially if its values are misaligned with human values.
			\item \textbf{Why A is wrong:} Speed is different from intelligence level
			\item \textbf{Why C is wrong:} Superintelligence is ABOVE human level
			\item \textbf{Why D is wrong:} Superintelligence would be general, not narrow
		\end{itemize}
		
		\textcolor{myred}{\textbf{EXAM TRAP:}} Superintelligence = \textcolor{myblue}{beyond human-level intelligence} - central to existential risk!
	\end{frame}
	
	\begin{frame}{Q82: AI Alignment - Tutorial}
		\fontsize{6.5}{7.5}\selectfont
		\textbf{QUESTION: What is AI alignment?}
		
		\vspace{0.1cm}
		\textbf{OPTIONS:}
		\begin{enumerate}
			\item Aligning AI systems with human values
			\item Aligning AI systems with maximum efficiency
			\item Aligning AI systems with profit maximization
			\item Aligning AI systems with speed
		\end{enumerate}
		
		\vspace{0.1cm}
		\textcolor{myblue}{\textbf{ANSWER: A}}
		
		\vspace{0.1cm}
		\textbf{DETAILED EXPLANATION:}
		\begin{itemize}
			\item \textbf{Why A is correct:} AI alignment is the challenge of ensuring that AI systems' means and goals are aligned with human values. Misaligned AI could cause catastrophic harm even with good intentions.
			\item \textbf{Why B is wrong:} Efficiency is not alignment\item \textbf{Why C is wrong:} Profit is not alignment
			\item \textbf{Why D is wrong:} Speed is not alignment
		\end{itemize}
		
		\textcolor{myred}{\textbf{EXAM TRAP:}} AI alignment = \textcolor{myblue}{matching AI to human values} - critical for safe AI!
	\end{frame}
	
	\begin{frame}{Q83: AI Arms Race - Tutorial}
		\fontsize{6.5}{7.5}\selectfont
		\textbf{QUESTION: What is the AI arms race concern?}
		
		\vspace{0.1cm}
		\textbf{OPTIONS:}
		\begin{enumerate}
			\item Competing nations developing AI for peaceful purposes
			\item Risks from different stakeholders creating powerful AI without adequate safety measures
			\item A race to create the fastest AI
			\item A race to create the most profitable AI
		\end{enumerate}
		
		\vspace{0.1cm}
		\textcolor{myblue}{\textbf{ANSWER: B}}
		
		\vspace{0.1cm}
		\textbf{DETAILED EXPLANATION:}
		\begin{itemize}
			\item \textbf{Why B is correct:} The AI arms race concern is that competing stakeholders may create powerful AI systems without adequate safety measures, potentially leading to catastrophic conflicts.
			\item \textbf{Why A is wrong:} The concern is about UNSAFE development
			\item \textbf{Why C is wrong:} Speed is not the primary concern
			\item \textbf{Why D is wrong:} Profit is not the primary concern
		\end{itemize}
		
		\textcolor{myred}{\textbf{EXAM TRAP:}} AI arms race = \textcolor{myblue}{unsafe AI development competition} - could lead to catastrophic outcomes!
	\end{frame}
	
	\begin{frame}{Q84: Economic Risks from AI - Tutorial}
		\fontsize{6.5}{7.5}\selectfont
		\textbf{QUESTION: What is a major economic risk from AI?}
		
		\vspace{0.1cm}
		\textbf{OPTIONS:}
		\begin{enumerate}
			\item AI always creates more jobs
			\item AI never affects employment
			\item Displacement of jobs by AI and widening inequality
			\item AI has no economic impact
		\end{enumerate}
		
		\vspace{0.1cm}
		\textcolor{myblue}{\textbf{ANSWER: C}}
		
		\vspace{0.1cm}
		\textbf{DETAILED EXPLANATION:}
		\begin{itemize}
			\item \textbf{Why C is correct:} AI could displace many jobs (high-skill jobs like accounting are also at risk), creating a skills gap and widening inequality. Predictions vary but Goldman Sachs suggests 1/4 of jobs could be automated.
			\item \textbf{Why A is wrong:} AI may DESTROY some jobs
			\item \textbf{Why B is wrong:} AI DOES affect employment
			\item \textbf{Why D is wrong:} AI HAS significant economic impact
		\end{itemize}
		
		\textcolor{myred}{\textbf{EXAM TRAP:}} Economic risks = \textcolor{myblue}{job displacement and inequality} - requires workforce re-skilling!
	\end{frame}
	
	\begin{frame}{Q85: Global Inequality from AI - Tutorial}
		\fontsize{6.5}{7.5}\selectfont
		\textbf{QUESTION: How could AI exacerbate global inequality?}
		
		\vspace{0.1cm}
		\textbf{OPTIONS:}
		\begin{enumerate}
			\item AI is evenly distributed globally
			\item Nations at the center of AI development leap ahead, leaving others behind
			\item AI has no effect on global inequality
			\item AI reduces inequality automatically
		\end{enumerate}
		
		\vspace{0.1cm}
		\textcolor{myblue}{\textbf{ANSWER: B}}
		
		\vspace{0.1cm}
		\textbf{DETAILED EXPLANATION:}
		\begin{itemize}
			\item \textbf{Why B is correct:} Most AI firms are based in the US, and technological adoption varies widely. Nations at the center of AI development could leap ahead in economic growth and influence, leaving others behind.
			\item \textbf{Why A is wrong:} AI is NOT evenly distributed
			\item \textbf{Why C is wrong:} AI CAN exacerbate inequality
			\item \textbf{Why D is wrong:} AI does NOT automatically reduce inequality
		\end{itemize}
		
		\textcolor{myred}{\textbf{EXAM TRAP:}} AI could \textcolor{myblue}{widen global inequality} - creates geopolitical tensions!
	\end{frame}
	
	% ============================================================
	% SECTION 9: GENAI RISKS - QUESTIONS 86-100
	% ============================================================
	
	\begin{frame}{Q86: GenAI Adoption Statistics - Tutorial}
		\fontsize{6.5}{7.5}\selectfont
		\textbf{QUESTION: According to McKinsey (2025), what percentage of organizations use GenAI?}
		
		\vspace{0.1cm}
		\textbf{OPTIONS:}
		\begin{enumerate}
			\item 33\%
			\item 71\%
			\item 50\%
			\item 95\%
		\end{enumerate}
		
		\vspace{0.1cm}
		\textcolor{myblue}{\textbf{ANSWER: B}}
		
		\vspace{0.1cm}
		\textbf{DETAILED EXPLANATION:}
		\begin{itemize}
			\item \textbf{Why B is correct:} McKinsey 2025 found that 71\% of organizations regularly use GenAI, up from 33\% in 2023. Adoption is particularly high in digitized industries like banking (58\% embrace GenAI).
			\item \textbf{Why A is wrong:} This was the 2023 figure
			\item \textbf{Why C is wrong:} The figure is higher than 50\%
			\item \textbf{Why D is wrong:} This is the BCG figure for professional services
		\end{itemize}
		
		\textcolor{myred}{\textbf{EXAM TRAP:}} GenAI adoption = \textcolor{myblue}{71\% of organizations} (2025 McKinsey) - rapid growth!
	\end{frame}
	
	\begin{frame}{Q87: Hallucination Definition - Tutorial}
		\fontsize{6.5}{7.5}\selectfont
		\textbf{QUESTION: What is hallucination in GenAI?}
		
		\vspace{0.1cm}
		\textbf{OPTIONS:}
		\begin{enumerate}
			\item Always correct outputs
			\item Factually incorrect, misleading, or fabricated outputs that appear credible
			\item Outputs that are never credible
			\item Outputs that are always correct
		\end{enumerate}
		
		\vspace{0.1cm}
		\textcolor{myblue}{\textbf{ANSWER: B}}
		
		\vspace{0.1cm}
		\textbf{DETAILED EXPLANATION:}
		\begin{itemize}
			\item \textbf{Why B is correct:} Hallucinations are factually incorrect or misleading outputs that appear credible to non-experts. They are particularly problematic in legal, financial, and healthcare contexts.
			\item \textbf{Why A is wrong:} Hallucinations are INCORRECT
			\item \textbf{Why C is wrong:} Hallucinations APPEAR credible
			\item \textbf{Why D is wrong:} Hallucinations are NOT correct
		\end{itemize}
		
		\textcolor{myred}{\textbf{EXAM TRAP:}} Hallucination = \textcolor{myblue}{incorrect but credible outputs} - hard for non-experts to detect!
	\end{frame}
	
	\begin{frame}{Q88: Air Canada Chatbot Case - Tutorial}
		\fontsize{6.5}{7.5}\selectfont
		\textbf{QUESTION: What happened in the Air Canada chatbot case?}
		
		\vspace{0.1cm}
		\textbf{OPTIONS:}
		\begin{enumerate}
			\item The chatbot worked perfectly
			\item The chatbot offered a discount that shouldn't have been available, and the company was held responsible
			\item The chatbot was a separate legal entity
			\item The company won the case
		\end{enumerate}
		
		\vspace{0.1cm}
		\textcolor{myblue}{\textbf{ANSWER: B}}
		
		\vspace{0.1cm}
		\textbf{DETAILED EXPLANATION:}
		\begin{itemize}
			\item \textbf{Why B is correct:} The chatbot offered a discount that shouldn't have been available. Air Canada argued the chatbot was a "separate legal entity," but the tribunal rejected this, ruling the company was responsible.
			\item \textbf{Why A is wrong:} The chatbot made an error
			\item \textbf{Why C is wrong:} The argument was REJECTED
			\item \textbf{Why D is wrong:} Air Canada LOST the case
		\end{itemize}
		
		\textcolor{myred}{\textbf{EXAM TRAP:}} Companies are \textcolor{myblue}{liable for chatbot outputs} - can't claim ignorance!
	\end{frame}
	
	\begin{frame}{Q89: Unpredictability in GenAI - Tutorial}
		\fontsize{6.5}{7.5}\selectfont
		\textbf{QUESTION: Why is unpredictability a risk in GenAI?}
		
		\vspace{0.1cm}
		\textbf{OPTIONS:}
		\begin{enumerate}
			\item GenAI outputs are always predictable
			\item Unpredictability undermines standardization across users and time
			\item Unpredictability is always beneficial
			\item Unpredictability never affects compliance
		\end{enumerate}
		
		\vspace{0.1cm}
		\textcolor{myblue}{\textbf{ANSWER: B}}
		
		\vspace{0.1cm}
		\textbf{DETAILED EXPLANATION:}
		\begin{itemize}
			\item \textbf{Why B is correct:} GenAI systems are unpredictable in their outputs, which undermines efforts to standardize outputs across users and time. This poses challenges for auditability and compliance.
			\item \textbf{Why A is wrong:} GenAI outputs are NOT always predictable
			\item \textbf{Why C is wrong:} Unpredictability is generally problematic
			\item \textbf{Why D is wrong:} Unpredictability CAN affect compliance
		\end{itemize}
		
		\textcolor{myred}{\textbf{EXAM TRAP:}} Unpredictability undermines \textcolor{myblue}{standardization and auditability} - compliance risk!
	\end{frame}
	
	\begin{frame}{Q90: Prompt Injection - Tutorial}
		\fontsize{6.5}{7.5}\selectfont
		\textbf{QUESTION: What is prompt injection?}
		
		\vspace{0.1cm}
		\textbf{OPTIONS:}
		\begin{enumerate}
			\item A technique to improve outputs
			\item Hidden instructions within user inputs to cause unintended behavior
			\item A method to increase accuracy
			\item A technique to remove bias
		\end{enumerate}
		
		\vspace{0.1cm}
		\textcolor{myblue}{\textbf{ANSWER: B}}
		
		\vspace{0.1cm}
		\textbf{DETAILED EXPLANATION:}
		\begin{itemize}
			\item \textbf{Why B is correct:} Prompt injection involves hidden instructions in user inputs designed to cause unintended behavior, manipulating the AI system's response.
			\item \textbf{Why A is wrong:} Prompt injection causes PROBLEMS
			\item \textbf{Why C is wrong:} Prompt injection is about manipulation
			\item \textbf{Why D is wrong:} Prompt injection doesn't remove bias
		\end{itemize}
		
		\textcolor{myred}{\textbf{EXAM TRAP:}} Prompt injection = \textcolor{myblue}{hidden instructions to cause unintended behavior} - security risk!
	\end{frame}
	
	\begin{frame}{Q91: Jailbreaking Definition - Tutorial}
		\fontsize{6.5}{7.5}\selectfont
		\textbf{QUESTION: What is jailbreaking in the context of GenAI?}
		
		\vspace{0.1cm}
		\textbf{OPTIONS:}
		\begin{enumerate}
			\item Making the model faster
			\item Bypassing a model's safety mechanisms through crafted prompts
			\item A method to improve outputs
			\item A technique to increase accuracy
		\end{enumerate}
		
		\vspace{0.1cm}
		\textcolor{myblue}{\textbf{ANSWER: B}}
		
		\vspace{0.1cm}
		\textbf{DETAILED EXPLANATION:}
		\begin{itemize}
			\item \textbf{Why B is correct:} Jailbreaking involves crafted prompts that bypass a model's safety mechanisms, overriding safety or policy constraints to cause unintended or unsafe behavior.
			\item \textbf{Why A is wrong:} Speed is not the purpose
			\item \textbf{Why C is wrong:} Jailbreaking is about bypassing safety
			\item \textbf{Why D is wrong:} Jailbreaking does NOT increase accuracy
		\end{itemize}
		
		\textcolor{myred}{\textbf{EXAM TRAP:}} Jailbreaking = \textcolor{myblue}{bypassing safety mechanisms} - security and safety risk!
	\end{frame}
	
	\begin{frame}{Q92: Data Privacy Violations in GenAI - Tutorial}
		\fontsize{6.5}{7.5}\selectfont
		\textbf{QUESTION: How can GenAI cause data privacy violations?}
		
		\vspace{0.1cm}
		\textbf{OPTIONS:}
		\begin{enumerate}
			\item GenAI never accesses personal data
			\item Prompts containing PII may be inadvertently stored, shared, or used for training
			\item Data privacy is not a concern
			\item GenAI always protects privacy
		\end{enumerate}
		
		\vspace{0.1cm}
		\textcolor{myblue}{\textbf{ANSWER: B}}
		
		\vspace{0.1cm}
		\textbf{DETAILED EXPLANATION:}
		\begin{itemize}
			\item \textbf{Why B is correct:} User prompts may contain personally identifiable information (PII) that can be inadvertently stored, shared, or used for model training. Researchers have extracted emails and phone numbers from LLMs.
			\item \textbf{Why A is wrong:} GenAI CAN access personal data
			\item \textbf{Why C is wrong:} Data privacy IS a concern
			\item \textbf{Why D is wrong:} GenAI does NOT always protect privacy
		\end{itemize}
		
		\textcolor{myred}{\textbf{EXAM TRAP:}} GenAI prompts may \textcolor{myblue}{inadvertently expose PII} - privacy risk!
	\end{frame}
	
	\begin{frame}{Q93: Copyright Issues in GenAI - Tutorial}
		\fontsize{6.5}{7.5}\selectfont
		\textbf{QUESTION: What copyright concern exists with GenAI?}
		
		\vspace{0.1cm}
		\textbf{OPTIONS:}
		\begin{enumerate}
			\item GenAI never uses copyrighted material
			\item GenAI models are trained on web-scraped content including copyrighted material
			\item Copyright laws don't apply to AI
			\item All GenAI content is public domain
		\end{enumerate}
		
		\vspace{0.1cm}
		\textcolor{myblue}{\textbf{ANSWER: B}}
		
		\vspace{0.1cm}
		\textbf{DETAILED EXPLANATION:}
		\begin{itemize}
			\item \textbf{Why B is correct:} GenAI models are trained on vast web-scraped content that includes copyrighted books, images, and articles. This has created legal proceedings (e.g., Getty Images vs. Stability AI).
			\item \textbf{Why A is wrong:} GenAI DOES use copyrighted material
			\item \textbf{Why C is wrong:} Copyright laws DO apply
			\item \textbf{Why D is wrong:} GenAI content is NOT public domain
		\end{itemize}
		
		\textcolor{myred}{\textbf{EXAM TRAP:}} GenAI training may \textcolor{myblue}{infringe copyright} - legal and reputational risk!
	\end{frame}
	
	\begin{frame}{Q94: Responsibility Gaps - Tutorial}
		\fontsize{6.5}{7.5}\selectfont
		\textbf{QUESTION: What are responsibility gaps in GenAI adoption?}
		
		\vspace{0.1cm}
		\textbf{OPTIONS:}
		\begin{enumerate}
			\item Clear lines of accountability always exist
			\item GenAI tools may be adopted across teams without clear responsibility for oversight
			\item Responsibility is never an issue
			\item All employees are always responsible
		\end{enumerate}
		
		\vspace{0.1cm}
		\textcolor{myblue}{\textbf{ANSWER: B}}
		
		\vspace{0.1cm}
		\textbf{DETAILED EXPLANATION:}
		\begin{itemize}
			\item \textbf{Why B is correct:} GenAI tools are intuitive and adopted across teams (including non-technical teams like HR), often without clear responsibility for oversight, quality assurance, and risk mitigation.
			\item \textbf{Why A is wrong:} Responsibility gaps ARE a problem
			\item \textbf{Why C is wrong:} Responsibility IS an issue
			\item \textbf{Why D is wrong:} Not all employees have clear responsibility
		\end{itemize}
		
		\textcolor{myred}{\textbf{EXAM TRAP:}} GenAI adoption creates \textcolor{myblue}{responsibility gaps} across teams - need clear accountability!
	\end{frame}
	
	\begin{frame}{Q95: Misaligned Incentives - Tutorial}
		\fontsize{6.5}{7.5}\selectfont
		\textbf{QUESTION: What misalignment exists in GenAI adoption?}
		
		\vspace{0.1cm}
		\textbf{OPTIONS:}
		\begin{enumerate}
			\item Companies always prioritize safety over speed
			\item Pressure to innovate can outweigh cautionary approaches
			\item Regulation always prevents problems
			\item Incentives are always aligned with safety
		\end{enumerate}
		
		\vspace{0.1cm}
		\textcolor{myblue}{\textbf{ANSWER: B}}
		
		\vspace{0.1cm}
		\textbf{DETAILED EXPLANATION:}
		\begin{itemize}
			\item \textbf{Why B is correct:} Pressure to cut costs and improve productivity can outweigh safety considerations. Teams may be rewarded for cost savings through automation, not for flagging risks or ensuring safe deployment.
			\item \textbf{Why A is wrong:} Safety is NOT always prioritized
			\item \textbf{Why C is wrong:} Regulation is not always effective
			\item \textbf{Why D is wrong:} Incentives are NOT always aligned
		\end{itemize}
		
		\textcolor{myred}{\textbf{EXAM TRAP:}} Innovation pressure can \textcolor{myblue}{outweigh safety considerations} - misaligned incentives!
	\end{frame}
	
	\begin{frame}{Q96: Lack of Training Risk - Tutorial}
		\fontsize{6.5}{7.5}\selectfont
		\textbf{QUESTION: What risk arises from lack of GenAI training?}
		
		\vspace{0.1cm}
		\textbf{OPTIONS:}
		\begin{enumerate}
			\item Employees always use GenAI correctly
			\item Employees may use GenAI without understanding acceptable use or data protection obligations
			\item Training is never needed
			\item GenAI always protects data
		\end{enumerate}
		
		\vspace{0.1cm}
		\textcolor{myblue}{\textbf{ANSWER: B}}
		
		\vspace{0.1cm}
		\textbf{DETAILED EXPLANATION:}
		\begin{itemize}
			\item \textbf{Why B is correct:} LLM chatbots are intuitive, so employees integrate them into workflows without guidance. 4.2\% of workers pasted sensitive data into ChatGPT, leading companies like JPMorgan to ban it.
			\item \textbf{Why A is wrong:} Employees may NOT use GenAI correctly
			\item \textbf{Why C is wrong:} Training IS needed
			\item \textbf{Why D is wrong:} GenAI does NOT always protect data
		\end{itemize}
		
		\textcolor{myred}{\textbf{EXAM TRAP:}} Lack of training leads to \textcolor{myblue}{unintentional data exposure} - significant risk!
	\end{frame}
	
	\begin{frame}{Q97: Vendor Lock-In - Tutorial}
		\fontsize{6.5}{7.5}\selectfont
		\textbf{QUESTION: What is vendor lock-in in GenAI?}
		
		\vspace{0.1cm}
		\textbf{OPTIONS:}
		\begin{enumerate}
			\item Complete freedom to switch vendors
			\item Dependence on external providers, limiting auditability and leverage
			\item Vendors never change pricing
			\item All providers are identical
		\end{enumerate}
		
		\vspace{0.1cm}
		\textcolor{myblue}{\textbf{ANSWER: B}}
		
		\vspace{0.1cm}
		\textbf{DETAILED EXPLANATION:}
		\begin{itemize}
			\item \textbf{Why B is correct:} Most organizations rely on external providers from a small set of dominant suppliers, creating dependencies and power asymmetries. Providers may change pricing, models, or usage restrictions without warning.
			\item \textbf{Why A is wrong:} Switching is often difficult
			\item \textbf{Why C is wrong:} Vendors CAN change terms
			\item \textbf{Why D is wrong:} Providers differ significantly
		\end{itemize}
		
		\textcolor{myred}{\textbf{EXAM TRAP:}} Vendor lock-in = \textcolor{myblue}{dependence limiting control} - strategic risk!
	\end{frame}
	
	\begin{frame}{Q98: Synthetic Content Risks - Tutorial}
		\fontsize{6.5}{7.5}\selectfont
		\textbf{QUESTION: What risk does synthetic content pose?}
		
		\vspace{0.1cm}
		\textbf{OPTIONS:}
		\begin{enumerate}
			\item Synthetic content is always beneficial
			\item High-quality misleading content can be created at scale, damaging trust and institutions
			\item Synthetic content is easy to identify
			\item Synthetic content never affects elections
		\end{enumerate}
		
		\vspace{0.1cm}
		\textcolor{myblue}{\textbf{ANSWER: B}}
		
		\vspace{0.1cm}
		\textbf{DETAILED EXPLANATION:}
		\begin{itemize}
			\item \textbf{Why B is correct:} GenAI lowers the cost of creating high-quality misleading content (fake news, health misinformation, political content), which can damage trust in information and democratic institutions.
			\item \textbf{Why A is wrong:} Synthetic content CAN be harmful
			\item \textbf{Why C is wrong:} Synthetic content is often difficult to identify
			\item \textbf{Why D is wrong:} Synthetic content CAN affect elections
		\end{itemize}
		
		\textcolor{myred}{\textbf{EXAM TRAP:}} Synthetic content can \textcolor{myblue}{damage trust and institutions} - societal risk!
	\end{frame}
	
	\begin{frame}{Q99: Deepfakes Risks - Tutorial}
		\fontsize{6.5}{7.5}\selectfont
		\textbf{QUESTION: What risk do deepfakes pose?}
		
		\vspace{0.1cm}
		\textbf{OPTIONS:}
		\begin{enumerate}
			\item Deepfakes are always harmless
			\item Audio or video deepfakes can convincingly mimic individuals, causing fraud, defamation, or political manipulation
			\item Deepfakes are easy to detect
			\item Deepfakes never cause harm
		\end{enumerate}
		
		\vspace{0.1cm}
		\textcolor{myblue}{\textbf{ANSWER: B}}
		
		\vspace{0.1cm}
		\textbf{DETAILED EXPLANATION:}
		\begin{itemize}
			\item \textbf{Why B is correct:} Deepfakes can convincingly mimic public figures or private individuals, causing fraud (CEO persuaded to transfer €220,000), defamation, and political manipulation. Even low-effort impersonations can go viral.
			\item \textbf{Why A is wrong:} Deepfakes CAN cause harm
			\item \textbf{Why C is wrong:} Deepfakes are often difficult to detect
			\item \textbf{Why D is wrong:} Deepfakes CAN cause harm
		\end{itemize}
		
		\textcolor{myred}{\textbf{EXAM TRAP:}} Deepfakes can \textcolor{myblue}{cause fraud, defamation, and manipulation} - serious societal risk!
	\end{frame}
	
	\begin{frame}{Q100: GenAI Risk Management Summary - Tutorial}
		\fontsize{6.5}{7.5}\selectfont
		\textbf{QUESTION: What is the most comprehensive summary of GenAI risk management?}
		
		\vspace{0.1cm}
		\textbf{OPTIONS:}
		\begin{enumerate}
			\item GenAI risks are purely technical
			\item GenAI risks include functional, data-related, organizational, strategic, and societal risks, requiring integrated oversight and governance
			\item GenAI risks are not significant
			\item Only hallucinations matter
		\end{enumerate}
		
		\vspace{0.1cm}
		\textcolor{myblue}{\textbf{ANSWER: B}}
		
		\vspace{0.1cm}
		\textbf{DETAILED EXPLANATION:}
		\begin{itemize}
			\item \textbf{Why B is correct:} GenAI risks span multiple categories: functional (hallucinations, unpredictability), data-related (privacy, copyright), organizational (responsibility gaps, misaligned incentives), strategic (vendor lock-in, regulatory unpredictability), and societal (disinformation, deepfakes).
			\item \textbf{Why A is wrong:} Risks extend beyond technical
			\item \textbf{Why C is wrong:} GenAI risks ARE significant
			\item \textbf{Why D is wrong:} Multiple risk categories exist
		\end{itemize}
		
		\textcolor{myred}{\textbf{EXAM SUCCESS:}} Understand \textcolor{myblue}{all risk categories} and their interconnections for comprehensive risk management!
	\end{frame}
	
\end{document}